\chapter{Fields, centers and reconstruction maps}
\label{cbc:chap:comparisons}

A collision coefficient has endomorphisms before it has a universal
interpretation.  The preceding chapters constructed coefficients,
comparison maps and actions on specified carriers.  We now compute
which operations commute with all morphisms in the categories those
coefficients generate.  The first example has only two objects: the
regular normal line and its punctured localization.  Their derived
center is a formal power-series algebra with a logarithmic odd
operation.  The arrow from the regular line to the punctured line
forces that logarithmic condition.

This calculation leads to two different extensions.  Several punctures
produce completed path algebras and transport on their centers.
The Heisenberg field algebra instead gives a tower of cochain actions
whose products retain a relative collision residue.  A universal action
must also account for maps and homotopies between actors.  Those maps
will supply the final obstruction: two parallel maps of pointed actions
remain nonhomotopic after every receiving pointed factorizing action.

\section{The regular and punctured normal line}
\label{cbc:rc:subsec:endomorphisms}

Let \(k\) be a field of characteristic zero.  Put
\[
 R=k[u],\qquad L=k[u,u^{-1}],\qquad
 D=k\langle u,\partial\rangle/(\partial u-u\partial-1).
\]
We regard \(R,L\) as left \(D\)-modules with their derivative connections.
Example~\ref{cbc:hc:binary} identifies the binary primitive coefficient
with \(R\oplus L\).  Its normal derivative is part of the coefficient.
We therefore take
\(\mathcal C_2=\operatorname{thick}_D(R\oplus L)\), the full dg category
generated by these modules under shifts, finite cones and direct summands.

For \(v\in D\), write \(r_v\) for right multiplication by \(v\), a map
of free left \(D\)-modules.  Use the resolutions
\begin{equation}\label{cbc:rc:eq:resolutions}
 P_0=[D\xrightarrow{r_\partial}D]\longrightarrow R,
 \qquad
 P_1=[D\xrightarrow{r_{u\partial+1}}D]\longrightarrow L
\end{equation}
in degrees \(-1,0\), augmented by \(1\) and \(u^{-1}\), respectively.
The Weyl algebra is a domain, so both differentials are injective.

The first cokernel is \(R\).  In the second, the relation
\(u\partial=-1\) on the cyclic generator reduces mixed normally ordered
monomials to pure powers of \(u\) or \(\partial\).
Their images are the distinct Laurent monomials
\(u^{m-1}\), \(m\geq0\), and \((-1)^n n!u^{-n-1}\), \(n\geq1\).
This proves the second cokernel identification, including injectivity of
its map to \(L\).

Applying \(\operatorname{Hom}_D(P_i,-)\), and changing the sign in
cochain degree one if necessary, gives
\[
\begin{array}{c|cc}
 &R&L\\\hline
 R &[R\xrightarrow{\partial}R]&[L\xrightarrow{\partial}L]\\
 L &[R\xrightarrow{u\partial+1}R]&[L\xrightarrow{u\partial+1}L].
\end{array}
\]
Polynomial integration, Laurent integration away from the residue, and
the formula \((u\partial+1)u^m=(m+1)u^m\) therefore yield
\begin{equation}\label{cbc:rc:eq:corners}
\begin{array}{c|cc}
 \operatorname{RHom}_D(-,-)&R&L\\\hline
 R &k&k\oplus k[-1]\\
 L &0&k\oplus k[-1].
\end{array}
\end{equation}
Here \(k[-1]\) lies in cohomological degree one.

To retain composition, define a graded path algebra \(T\) on objects
\(0,1\). Its generators are \(a:0\to1\) of degree zero and
\(\theta:1\to1\) of degree one, subject to \(\theta^2=0\).
Products are written in composition order.  Put \(b=\theta a\).
Thus
\[
 T^0=k e_0\oplus k e_1\oplus k a,\qquad
 T^1=k\theta\oplus k b,\qquad d_T=0.
\]

Apart from identities, the only nonzero product of basis elements is
\(\theta a=b\).
We use perfect right dg modules in \(\operatorname{Perf}(T)\).

\begin{theorem}\label{cbc:rc:thm:endomorphisms}
There is a strict quasi-isomorphism of dg algebras
\[
 T\longrightarrow\operatorname{End}_D(P_0\oplus P_1).
\]
Consequently \(\mathcal C_2\simeq\operatorname{Perf}(T)\), as dg categories.
\end{theorem}
\begin{proof}
Map \(e_i\) to the projections.  Map \(a\) to \(\alpha:P_0\to P_1\), with
\(\alpha^{-1}=1\), \(\alpha^0=r_u\).  Its chain equation is
\(\partial u=u\partial+1\).
Map \(\theta\) to the degree-one identity
\(\Theta:P_1^{-1}\to P_1^0\), and \(b\) to \(\Theta\alpha\).
These maps are closed and satisfy \(\Theta^2=0\) strictly.

After augmentation, \(\alpha\) represents the inclusion \(R\to L\).
\(\Theta\) and \(\Theta\alpha\) represent \(u^{-1}\) in the Euler and
derivative cokernels.  They give bases of all the corners in
\eqref{cbc:rc:eq:corners}.  The resulting algebra map is a quasi-isomorphism.
The functor \(\operatorname{Hom}_D(P_0\oplus P_1,-)\) and its tensor
adjoint identify the generators and their Hom complexes.  Exactness and
passage to direct summands extend this identification to their thick
closures, proving the last assertion.
\end{proof}

Restrict the auxiliary binary Heisenberg complex to
\((z_1,z_2)=(0,u)\), with its induced normal connection, and denote it
by \(Z^{\mathrm{nor}}_{\kappa,2}\).  At nonzero level,
Theorem~\ref{cbc:hc:polynomial-cohomology} gives the coefficient decomposition
\[
 Z^{\mathrm{nor}}_{\kappa,2}\simeq
 R\oplus(R\oplus L)[-1]\oplus L[-2].
\]
It generates the same dg category.  This statement concerns the
coefficient complex. The native transverse multiplication supplies
additional structure.

\section{The completed center of the normal pair}
\label{cbc:rc:subsec:centre}

For a dg category generated by \(T\), its categorical center is the
derived endomorphism complex of the diagonal \(T\)-bimodule, with its
Hochschild operations.  We compute it using the normalized relative
two-sided bar over \(S=k e_0\oplus k e_1\).
The algebra \(S\) is separable.  Removing the unit degeneracies from
the two-sided bar therefore leaves a resolution
\[
 T\otimes_S T^c_S((T/S)[1])\otimes_S T\longrightarrow T.
\]
Here \(T^c_S(W)=\bigoplus_{n\geq0}W^{\otimes_S n}\), with the
length-zero term \(S\), is the tensor coalgebra over \(S\).

Applying
\(\operatorname{Hom}_{T^e}(-,T)\) gives the \emph{product} over arities.
Denote this cochain algebra by \(C(T)\).
This construction and dg Morita invariance identify it with the center
of \(\mathcal C_2\). See \cite[\S5.4]{Keller2006DG}.

Fix the following signs.  An \(n\)-cochain of internal degree \(r\)
has total degree \(n+r\).  Its differential is \(D=(-1)^r b_H\), where
\begin{align*}
 (b_H f)(v_1,\ldots,v_{n+1})={}&
 (-1)^{r|v_1|}v_1f(v_2,\ldots,v_{n+1})\\
 &+\sum_{i=1}^{n}(-1)^i
 f(v_1,\ldots,v_iv_{i+1},\ldots,v_{n+1})\\
 &+(-1)^{n+1}f(v_1,\ldots,v_n)v_{n+1}.
\end{align*}
For an \(m\)-cochain \(g\) of internal degree \(s\), use
\begin{equation}\label{cbc:rc:eq:cup}
 (f\smile g)(v_1,\ldots,v_{n+m})=
 (-1)^{s(\sum_{i=1}^n|v_i|-n)}
 f(v_1,\ldots,v_n)g(v_{n+1},\ldots,v_{n+m}).
\end{equation}
These are the suspended-bar differential and product expressed on
unsuspended inputs. In particular \(D\) satisfies the total-degree
Leibniz rule.

\begin{theorem}\label{cbc:rc:thm:centre}
There is a continuous strict quasi-isomorphism of associative dg algebras
\begin{equation}\label{cbc:rc:eq:centre}
 k[[x]]\otimes\Lambda(\eta)\longrightarrow C(T),
 \qquad |x|=0,\quad |\eta|=1,
\end{equation}
with zero differential on the source.  The Gerstenhaber bracket on
cohomology satisfies
\[
 [\eta,x]=-x,\qquad [x,x]=[\eta,\eta]=0.
\]
\end{theorem}
\begin{proof}
Assign weight one to \(\theta,b\) and weight zero to \(e_i,a\).
For a cochain, let \(d\) be input weight minus output weight.
The differential preserves \(d\).

Every nonempty composable relative word is
\(\theta^{\otimes n}\), \(\theta^{\otimes(n-1)}\otimes a\), or
\(\theta^{\otimes(n-1)}\otimes b\).
The output must have the same endpoints as its word.  In weight
\(d\geq1\), this gives exactly the following ordered bases:
\[
\begin{array}{c|l}
0&\theta^{\otimes d}\mapsto e_1,\quad
  \theta^{\otimes(d-1)}\otimes b\mapsto a\\
1&\theta^{\otimes(d+1)}\mapsto\theta,\quad
  \theta^{\otimes d}\otimes a\mapsto a,\quad
  \theta^{\otimes d}\otimes b\mapsto b\\
2&\theta^{\otimes(d+1)}\otimes a\mapsto b.
\end{array}
\]
The differential matrices are
\begin{equation}\label{cbc:rc:eq:blocks}
 k^2\xrightarrow{M}k^3\xrightarrow{N}k,
 \qquad
 M=\begin{pmatrix}0&0\\-1&1\\-1&1\end{pmatrix},
 \qquad N=\begin{pmatrix}1&1&-1\end{pmatrix}.
\end{equation}
Indeed \(\theta a=b\) is the only possible nonzero interior face.
The two endpoint faces give the other entries.

In weight zero the degree-zero basis is \(e_0,e_1\), and the first
matrix is \(-M\).  Weight \(-1\) consists of the identity map
\(k\to k\) in degrees one and two. Its source is the nullary
\(\theta\), and its target sends \(a\) to \(b\).
No lower weight occurs.

For each \(d\geq0\), define representatives and projections by
\[
 i_0(c)=(c,c),\quad p_0(p,q)=p,\qquad
 i_1(c)=(c,0,c),\quad p_1(r,s,t)=t-s.
\]
A contracting homotopy off these representatives is
\[
 h_1(r,s,t)=
 \begin{cases}(0,s)&d>0,\\(0,-s)&d=0,\end{cases}
 \qquad h_2(z)=(z,0,0).
\]
Matrix multiplication gives \(Dh+hD=1-ip\) and
\(h^2=hi=ph=0\).  Use the inverse identity in weight \(-1\).
Every block is finite, and these maps preserve weight, so their product
contracts the full cochain complex continuously.

Define derivations \(X,Y\) by
\[
\begin{array}{c|ccccc}
 &e_0&e_1&\theta&a&b\\\hline
 X&0&0&e_1&0&a\\
 Y&0&0&\theta&0&b.
\end{array}
\]
Their internal degrees are \(-1,0\), and their total degrees are
\(0,1\).  The displayed product table proves the graded derivation
identities, hence \(DX=DY=0\).
Formula~\eqref{cbc:rc:eq:cup} gives
\(X\smile Y=Y\smile X\) and \(Y\smile Y=0\).

The powers \(X^{\smile d}\) and \(Y\smile X^{\smile d}\) are the
representatives \((1,1)\) and \((1,0,1)\) above.
Thus \(x\mapsto X\), \(\eta\mapsto Y\) defines the asserted
quasi-isomorphism.  Each formal coefficient has a specified arity, so
evaluation on a fixed word involves a finite sum.
The unary Gerstenhaber bracket is the graded commutator of derivations.
Here \(YX-XY=-X\).  The biderivation rule on cohomology gives the bracket
on all formal series.
\end{proof}

Thus \(\eta\) acts as the logarithmic vector field \(-x\partial_x\).
In particular
\[
 [\eta f,g]=-xf g',\qquad
 [\eta f,\eta g]=\eta x(gf'-fg').
\]
The full bar, with its braces, retains the higher center operations.
The associative quasi-isomorphism and the cohomology bracket above do
not assert formality as an \(E_2\)-algebra.
Using direct sums of cochains instead would replace \(k[[x]]\) by
\(k[x]\). That convention omits endomorphisms of the direct-sum bar.

The restriction maps explain the logarithmic condition.
The full subcategory on \(L\) has endomorphism algebra
\(E=\Lambda(\theta)\), and its normalized cochains are
\(C(E)=E[[x]]\), with zero differential.  Restricting relative words gives
\begin{equation}\label{cbc:rc:eq:restrictions}
 k[[x]]\otimes\Lambda(\eta)\longrightarrow E[[x]],
 \qquad x\longmapsto x,\quad\eta\longmapsto\theta x.
\end{equation}
Restriction to \(R\) sends \(x,\eta\) to zero and preserves constants.
These are restrictions along full dg subcategories.
The extra regular object forces the odd coefficient in \(E[[x]]\)
to vanish at the formal origin.

One can also see this from the relative bar triangle
\[
 C(T)\longrightarrow C(k)\oplus C(E)
 \longrightarrow\operatorname{RHom}_E(Ea,Ea).
\]
Split words into those in a diagonal corner and those crossing from
\(0\) to \(1\).  Crossing occurs once. The latter words are the shifted
module bar of \(Ea\).  Since \(Ea\) is free over \(E\), that bar contracts
to \(E\).

On the diagonal models the last arrow is
\((c,f)\mapsto c-f(0)\), a surjective map of complexes.
Its kernel is \(k\oplus xE[[x]]\). The identification
\(\eta=\theta x\) recovers \eqref{cbc:rc:eq:centre}.

\section{Marked deformations and the Kummer residue}
\label{cbc:rc:subsec:deformation}

First fix the two objects, or equivalently the inclusion
\(S=ke_0\oplus ke_1\subset T\).  For uncurved \(S\)-linear
deformations with zero nullary comparison term, retain only positive
arities.  Let \(C_{\geq1}(T)\) denote this subcomplex of \(C(T)\),
with the same product totalization.  Write \(C^k_{\geq1}(T)\) for
the positive-arity normalized Hochschild complex relative to \(k\)
instead of \(S\). It fixes the unit but permits the object
decomposition to move.

\begin{proposition}\label{cbc:rc:prop:positive-arity}
The \(S\)-relative complex \(C_{\geq1}(T)\) has cohomology
\[
 H^0=xk[[x]],\qquad
 H^1=k[[x]]\eta\oplus k\sigma,\qquad
 H^2=kq,
\]
and vanishes otherwise.  Here \(\sigma(a)=a\), \(\sigma(b)=b\),
and \(q(a)=b\). Their other generator values vanish.
For the complex relative to \(k\), the answer is
\[
 H^0=xk[[x]],\qquad
 H^1=k[[x]]\eta\oplus k\sigma\oplus kDa,\qquad
 H^2=kD\theta\oplus kDb,
\]
with no other cohomology.  The last two generators have total degree
two, and \(D\theta=q\).
\end{proposition}
\begin{proof}
The positive-weight blocks in \eqref{cbc:rc:eq:blocks} do not change.
Removing the two nullary identities in weight zero removes their
degree-zero cohomology and leaves their former boundary
\((0,1,1)=\sigma\) in degree one.
Removing the nullary \(\theta\) in weight \(-1\) leaves \(q\) in
degree two.  The enumeration in the preceding proof exhausts the
cochains, and gives the first formula.

For the second assertion, restore arity zero.  Its quotient from the
full \(k\)-relative complex is \(T\), with zero differential, giving
\[
 0\longrightarrow C^k_{\geq1}(T)\longrightarrow C^k(T)
   \longrightarrow T\longrightarrow0.
\]
Separability of \(S\) identifies the cohomology of \(C^k(T)\) with
that computed in Theorem~\ref{cbc:rc:thm:centre}.  Projection to arity
zero sends \(f(x)\) to \(f(0)1_T\) and sends
\(\eta k[[x]]\) to zero: all these latter representatives have
positive arity.

The cohomology exact sequence therefore contributes
\(T^0/k1_T\), with representatives \(e_0,a\), to degree one, and
all of \(T^1=k\theta\oplus kb\) to degree two.
The connecting maps are the Hochschild differentials of these
nullary elements.  In particular \((Db)(e_0)=b\), whereas
\((D\theta)(e_0)=0\) and \((D\theta)(a)=b\), so the two
degree-two classes are distinct.  The positive-arity representatives
\(\eta k[[x]]\) split the degree-one exact sequence.
\end{proof}

Put \(A=k[[s]]\).  Let \(T_s\) have the graded multiplication of
\(T\otimes_k A\), but differential
\begin{equation}\label{cbc:rc:eq:deformation}
 d_s a=s b,\qquad d_s\theta=d_s b=d_s e_i=0.
\end{equation}
It is a dg algebra: \(d_s=s[\theta,-]\) is a derivation of square zero.
Its graded \(A\)-module is free, whereas
\[
 H^0(T_s)=Ae_0\oplus Ae_1,\qquad
 H^1(T_s)=A\theta\oplus(A/(s))b.
\]

Even modulo \(s^2\), there is no uncurved \(A_\infty\)-map
\(T_s\to T\otimes A\) with \(F_0=0\) and \(F_1=1\pmod{s}\).
Its unary equation would be \(F_1d_s=0\), but on \(a\) the left side
is \(s b\ne0\pmod{s^2}\).
Higher components do not enter this equation.
The full Hochschild complex has \(D\theta=q\) and \(H^2C(T)=0\),
but its primitive is excluded from this comparison problem.
The example concerns graded flatness. Its cohomology is not flat over
\(A\).

The family has a realization by differential modules.  Define
\(D_A=A\langle u,\partial\rangle/(\partial u-u\partial-1)\),
\(R_A=A[u]\) with its derivative connection, and
\[
 L_s=A[u,u^{-1}],\qquad \partial f=f'-\frac{s}{u}f.
\]
These are polynomial and Laurent polynomial modules over \(A\).
No completion in \(u\) is taken.

\begin{theorem}\label{cbc:rc:thm:kummer}
There is a strict quasi-isomorphism
\[
 T_s\longrightarrow
 \operatorname{End}_{D_A}(P_{0,A}\oplus P_{1,s}),\qquad
 P_{1,s}=[D_A\xrightarrow{r_{u\partial+1+s}}D_A],
\]
where \(P_{0,A}=P_0\otimes_k A\), and \(P_{1,s}\to L_s\) is augmented
by \(u^{-1}\).  It induces
\[
 \operatorname{Perf}(T_s)\simeq
 \operatorname{thick}_{D_A}(R_A\oplus L_s).
\]
For \(K=k((s))\), the generic fiber satisfies
\[
 \operatorname{Perf}(T_s\otimes_A K)\simeq
 \operatorname{Perf}(K)\oplus\operatorname{Perf}(\Lambda_K(\theta)).
\]
Consequently \(T_s\) is not Morita equivalent over \(A\) to the
constant family \(T\otimes_k A\).
\end{theorem}
\begin{proof}
The ordered Weyl filtration proves that \(D_A\) is a domain.
In the cyclic quotient by \(u\partial+1+s\), mixed monomials reduce
to pure powers: commuting one \(\partial\) past \(u\) lowers both
exponents, using the cyclic Euler relation.
The pure powers act on \(u^{-1}\) as
\[
 u^m u^{-1}=u^{m-1},\qquad
 \partial^n u^{-1}=(-1)^n\prod_{j=1}^n(j+s)u^{-n-1}.
\]
The factors \(j+s\) are units of \(A\). These images form a Laurent
basis.  This proves exactness of the displayed resolution.

Keep the maps \(\alpha^{-1}=1\), \(\alpha^0=r_u\), \(\Theta\), and \(\Theta\alpha\)
from Theorem~\ref{cbc:rc:thm:endomorphisms}.  Now
\[
 d_{\operatorname{Hom}}\alpha
   =r_{u\partial+1+s}-r_{\partial u}=s\Theta\alpha.
\]
All other differentials and all multiplicative relations agree with
\eqref{cbc:rc:eq:deformation}, giving a strict dg algebra map.

Compute its four corners after augmenting the target resolutions.
For \(R_A\to L_s\), the derivative maps
\(u^m\) to \((m-s)u^{m-1}\).  All modes with \(m\ne0\) contract,
since \(m-s\) is a unit.  With the Hom differential sign, the remaining
complex is \([A\xrightarrow{s}A]\), represented by \(\alpha,\Theta\alpha\).
For \(L_s\to R_A\), the Euler differential acts by \(m+1+s\) on
\(u^m\), \(m\geq0\). Every coefficient is a unit, so this corner is
acyclic.

For \(L_s\to L_s\), it acts by \(m+1\), leaving the
\(m=-1\) mode in both degrees.  The remaining regular corner is \(A\)
in degree zero.  Thus the map is a quasi-isomorphism on every corner.
The same generator argument as before proves the dg equivalence.

Over \(K\), the two-term ideal spanned by \(a,b\) is contractible,
because its differential is multiplication by the unit \(s\).
The quotient map
\(T_s\otimes_A K\to K\times\Lambda_K(\theta)\) is therefore a
dg algebra quasi-isomorphism, proving the generic-fiber assertion.
Its degree-zero categorical center is \(K\times K[[x]]\).
Indeed the center of a product category is the product of its
centers, and the exterior corner has the cochain model
\(\Lambda_K(\theta)[[x]]\) already computed above.

Applying Theorem~\ref{cbc:rc:thm:centre} directly over the field \(K\)
gives \(HH^0(T\otimes_k K)=K[[x]]\).
The first ring has a nontrivial idempotent, whereas the second has
none.  Morita equivalence would preserve these center rings. An
\(A\)-linear Morita equivalence would remain one after extension to
\(K\).  Neither equivalence exists.  This argument computes both
generic centers directly and uses no base-change assertion for
completed Hochschild cochains.
\end{proof}

The special fiber has a horizontal inclusion \(R\to L\).
In the family, the ordinary inclusion of coefficient modules fails to
commute with the connections by \(-s/u\). The equation \(d_s a=s b\)
records this failure.  Thus the generic splitting is caused by the loss
of that horizontal arrow.  It makes no assertion that completed
Hochschild cochains commute with specialization in \(s\).

\section{Independent normal directions}
\label{cbc:rc:subsec:crossings}

Let \(R_d=k[u_1,\ldots,u_d]\), and let \(D_d\) be its Weyl algebra.
The normal-crossing coefficient generator is
\[
 G_d=\bigoplus_{I\subseteq\{1,\ldots,d\}}
 R_d\left[\left(\prod_{i\in I}u_i\right)^{-1}\right]
 =\bigotimes_{i=1}^d(k[u_i]\oplus k[u_i,u_i^{-1}]),
\]
where the tensor products are external over \(k\) and the empty product
is one.  Each summand has its derivative connection.

\begin{theorem}\label{cbc:rc:thm:normal-crossing}
The categorical center of \(\operatorname{thick}_{D_d}(G_d)\), with full
product-total Hochschild cochains, has the associative dg model
\[
 k[[x_1,\ldots,x_d]]\otimes\Lambda(\eta_1,\ldots,\eta_d),
 \qquad |x_i|=0,\quad |\eta_i|=1,
\]
with zero differential.  Its cohomology bracket satisfies
\[
 [\eta_i,x_j]=-\delta_{ij}x_j,\qquad
 [\eta_i,\eta_j]=[x_i,x_j]=0.
\]
\end{theorem}
\begin{proof}
Tensor the bounded free resolutions \eqref{cbc:rc:eq:resolutions}.
Their endomorphism algebra identifies with the tensor product of the
individual endomorphism algebras, with the Koszul signs.
The tensor of the strict maps in Theorem~\ref{cbc:rc:thm:endomorphisms}
is a quasi-isomorphism from \(T^{\otimes d}\).  It gives the required
dg Morita comparison.

Use the tensor product of the relative two-sided bars as a resolution
of its diagonal bimodule.  The tensor of their bar diagonals gives the
cup product.  The shuffle and Alexander--Whitney comparisons with the
bar of \(T^{\otimes d}\) identify the resulting derived endomorphism
algebras. They also identify insertions up to coherent homotopy.
This is a comparison of resolutions, not an equality of their terms.

At each multiweight, the cochain spaces are finite-dimensional and
bounded.  Hence tensor Hom identifies this model with the tensor of
the corresponding binary cochain blocks.
Their contractions tensor: for two factors the homotopy is
\(h_1\otimes1+i_1p_1\otimes h_2\), with the graded tensor sign.
Iteration gives all \(d\) factors.

Taking products over multiweights gives precisely the formal power
series in the statement.  The multiplicative maps of
Theorem~\ref{cbc:rc:thm:centre} tensor to an associative quasi-isomorphism
to this diagonal model.
The insertion comparison makes the bracket the sum of the factor
brackets.  Operations in distinct factors have zero mixed bracket.
\end{proof}

The theorem applies to a normal slice of two independent pair
collisions, such as \(12|34\), with constant tangential coefficients.
The exterior coefficient complex contains all four normal
localizations as shifted direct summands and generates this category.
A global cut open can carry additional tangential cohomology.
At a triple collision the three normal equations \(u,v,u-v\) are
dependent, so the preceding independent-direction hypothesis fails.
Its shape provides a different algebra.

\section{The shape of a triple collision}
\label{cbc:rc:subsec:shape}

On the distinct three-point configuration, use coordinates
\[
 z_3=z,\qquad z_1=z+u,\qquad z_2=z+ut.
\]
They identify it with
\(\mathbb A^1_z\times\mathbb G_{m,u}\times V\), where
\[
 V=\operatorname{Spec}k[t,t^{-1},(t-1)^{-1}].
\]
The scale is \(u\).  The two punctures of \(V\) retain the two finite
pair-collision channels within that scale.

Let
\[
 B=k\oplus k\epsilon_0\oplus k\epsilon_1,
 \qquad |\epsilon_i|=1,
 \qquad \epsilon_i\epsilon_j=0,
\]
with zero differential.  The map
\(\epsilon_0\mapsto dt/t\), \(\epsilon_1\mapsto dt/(t-1)\)
is a strict multiplicative quasi-isomorphism \(B\to\Omega^\bullet(V)\).
Indeed partial fractions integrate every polynomial term and every pole
of order at least two.  The two simple residues are independent, and
the horizontal functions are constants.  Products of the one-forms
vanish on the curve.

Thus \(B\) models the endomorphisms of the trivial differential module
on \(V\), and \(\Lambda(\rho)\otimes B\), with
\(\rho=du/u\), models its product with the scale direction.
The ordinary de Rham dimensions of the latter are \(1,3,2\).

\begin{theorem}\label{cbc:rc:thm:shape-centre}
Let \(F=k\langle\!\langle x,y\rangle\!\rangle\) be the completed free
associative algebra, with its word-length topology.
The normalized full Hochschild cochains of \(B\) identify with the dg
algebra
\begin{equation}\label{cbc:rc:eq:shape-centre}
 F\xrightarrow{d}F\oplus F,\qquad
 df=([x,f],[y,f]),
\end{equation}
in degrees zero and one.  The degree-one ideal has square zero, and
its left and right actions are componentwise multiplication by \(F\).
In particular
\[
 H^0C(B)=k,\qquad
 H^1C(B)=\frac{F\oplus F}
 {\{([x,f],[y,f]):f\in F\}},
 \qquad H^jC(B)=0\quad(j\ne0,1).
\]
The word-length \(m\) component of \(H^1\) has dimension
\(2\) for \(m=0\), \(4\) for \(m=1\), and
\(3\cdot2^{m-1}\) for \(m\geq2\).
\end{theorem}
\begin{proof}
Both suspended augmentation generators have degree zero.  A cochain
with scalar output is a formal series in their words, in total degree
zero.  Outputs in \(k\epsilon_0\oplus k\epsilon_1\) give two such
series, in total degree one.  All interior faces vanish, because the
augmentation ideal is square zero.  With the conventions of
\eqref{cbc:rc:eq:cup}, the endpoint faces give the two commutators in
\eqref{cbc:rc:eq:shape-centre}.  The cup product concatenates words and
multiplies outputs, giving exactly the asserted algebra.

A homogeneous word polynomial commuting with \(x\) is a polynomial
in \(x\).  To see this, the relation \(xf=fx\) first excludes every
word ending in \(y\). Cycling a final \(x\) to the beginning then
excludes every word containing \(y\).
Commutation with \(y\) leaves only constants.
Apply this argument in each length to obtain \(Z(F)=k\).

At output length \(m\), the differential comes from scalar words of
length \(m-1\).  Its rank is zero for \(m=1\), and
\(2^{m-1}\) for \(m\geq2\). There is no source for \(m=0\).
Subtracting these ranks from \(2\cdot2^m\) gives the dimensions.
The image is the product of its finite-dimensional length components,
so it is closed.  The quotient and cohomology are therefore compatible
with the stated completion.
\end{proof}

This computes the separated shape corner.  It does not identify the
full triple coefficient category with the category generated by that
one corner.

\section{Three-point resonance}
\label{cbc:rc:subsec:resonance}

For \(n\geq2\), \(c\in k\), and numerical charges \(q_1,\ldots,q_n\in k\), consider the rank-one
connection on the distinct configuration
\[
 \nabla_n=d+c\sum_{i<j}q_iq_j\,d\log(z_i-z_j).
\]
Normalize two positions and write every difference as \(u\) times
a shape function.  Its radial residue is
\[
 \lambda_n=c\sum_{i<j}q_iq_j.
\]
The radial differential on \(u^m\) is multiplication by
\(m+\lambda_n\), followed by \(du/u\).
If \(\lambda_n\notin\mathbb Z\subset k\), division by these scalars
contracts every Laurent mode.  The global algebraic de Rham complex
is then acyclic.

If \(\lambda_n\in\mathbb Z\), its radial cohomology
has representatives \(u^{-\lambda_n}\) and
\(u^{-\lambda_n}du/u\).  Tensoring with the shape complex computes
the remaining cohomology.  The translation coordinate contracts to
constants by polynomial integration.

\begin{proposition}\label{cbc:rc:prop:ternary-resonance}
Take \(q_1=q_2=q_3=1\) and \(c=1/3\).
Every two-point global algebraic de Rham complex is acyclic, but the
three-point complex has
\[
 H^0=0,\qquad
 H^1=k\left[u^{-1}\frac{dt}{t}\right],\qquad
 H^2=k\left[u^{-1}\frac{du}{u}\wedge\frac{dt}{t}\right],
\]
and no other cohomology.
\end{proposition}
\begin{proof}
Each pair has radial residue \(1/3\), while the triple has residue one.
In the scale and shape coordinates its normal connection is
\[
 d+\frac{du}{u}+\frac13\left(\frac{dt}{t}+\frac{dt}{t-1}\right).
\]
It remains to compute the shape operator
\[
 D_{\mathrm{sh}}f=f'+
 \left(\frac1{3t}+\frac1{3(t-1)}\right)f
\]
on \(k[t,t^{-1},(t-1)^{-1}]\).
Every \(r\) has a decomposition
\begin{equation}\label{cbc:rc:eq:shape-reduction}
 r=D_{\mathrm{sh}}f+\frac{c(r)}t.
\end{equation}
For a pole at \(p=0,1\) of order \(j\geq2\) and leading coefficient
\(a\), subtract
\[
 D_{\mathrm{sh}}\left(
 \frac{a(t-p)^{1-j}}{1/3-(j-1)}\right).
\]
This lowers the pole order at \(p\) and introduces at most a simple
pole at the other point.

After eliminating higher poles, remove a
polynomial leading term \(a t^j\) by subtracting
\(D_{\mathrm{sh}}(a t^{j+1}/(j+1+2/3))\).
The remainder is \(A_0/t+A_1/(t-1)\).  Subtracting
\(D_{\mathrm{sh}}(3A_1)\) leaves \((A_0-A_1)/t\).
All denominators used are nonzero, and every reduction terminates.
For instance,
\[
 \frac1{t^2(t-1)^2}=
 D_{\mathrm{sh}}\left(\frac{3-6t^2}{2t(t-1)}\right)+\frac2t.
\]

There is no nonzero rational horizontal function: its order at zero
would have to be \(-1/3\).
Also \(1/t\) is not in the image of \(D_{\mathrm{sh}}\).
A pole of order \(m\geq1\) of a proposed primitive at either puncture
would produce a higher pole with coefficient \(-m+1/3\ne0\).
The primitive must therefore be polynomial.  Degree \(j\geq1\)
produces a nonzero polynomial part with coefficient \(j+2/3\), so the
primitive must be constant.  A constant produces equal coefficients
at the two simple poles, and cannot produce \(1/t\).

The decomposition \eqref{cbc:rc:eq:shape-reduction} thus defines a
contraction of the shape complex onto \(k\,dt/t\) in degree one:
its homotopy sends \(r\,dt\) to the unique \(f\).
The radial homotopy sends \(u^mdu/u\) to \(u^m/(m+1)\) for
\(m\ne-1\), and sends the remaining mode to zero.
Tensor these two contractions with the graded tensor signs.
The retained radial modes \(u^{-1}\), \(u^{-1}du/u\) give exactly the
two classes in the statement.
\end{proof}

The dependence on the three residues can be determined without
restricting them to products of charges. Let $k$ be a field of
characteristic zero, let $a_{12},a_{13},a_{23}\in k$, and set
$\lambda=a_{12}+a_{13}+a_{23}$. On the ordered configuration of
three points of $\mathbb A^1_k$, take the rank-one connection
\[
 \nabla=d+\sum_{i<j}a_{ij}\,d\log(z_i-z_j).
\]
All functions and differential forms below are algebraic, with
finite Laurent support.

\begin{proposition}\label{cbc:rc:prop:three-residues}
The dimensions of the global de Rham cohomology are
\[
\begin{array}{c|ccc}
 &\dim_k H^0&\dim_k H^1&\dim_k H^2\\ \hline
 \lambda\notin\mathbb Z&0&0&0\\
 \lambda\in\mathbb Z,\ (a_{12},a_{13},a_{23})\notin\mathbb Z^3
     &0&1&1\\
 (a_{12},a_{13},a_{23})\in\mathbb Z^3&1&3&2
\end{array}
\]
Here $\mathbb Z$ denotes its image in $k$. All other cohomology
groups vanish.
\end{proposition}

\begin{proof}
Put $z_1=z+u$, $z_2=z+ut$, $z_3=z$. The coordinate ring is
$k[z,u,u^{-1},t,t^{-1},(t-1)^{-1}]$, and the connection becomes
\[
 d+\lambda\frac{du}{u}
   +a_{23}\frac{dt}{t}+a_{12}\frac{dt}{t-1}.
\]
Polynomial integration contracts the translation variable. The
scale differential on $u^m$ is multiplication by $m+\lambda$.
It is contractible unless $\lambda$ is integral, in which case
it has one class in each of degrees zero and one.

Write $a=a_{23}$, $b=a_{12}$. The shape differential is
$D f=(f'+(a/t+b/(t-1))f)\,dt$ on
$B=k[t,t^{-1},(t-1)^{-1}]$. Choose an integer $M\geq1$ such
that none of $a-j$, $b-j$, or $j+a+b$ vanishes for an integer
$j>M$. Let $V_M\subset B$ consist of functions with poles of
order at most $M$ at zero and one and polynomial part of degree
at most $M$. Let $W_M\subset B\,dt$ consist of forms with pole
orders at most $M+1$ and polynomial part of degree at most $M-1$.
Then $D(V_M)\subset W_M$, and
$\dim V_M=3M+1$, $\dim W_M=3M+2$.

This finite subcomplex computes the shape cohomology. Indeed,
a pole of order $j\geq M+2$ is removed by the derivative of
a function of order $j-1$, since its leading coefficient is
$a-(j-1)$ or $b-(j-1)$. Such a subtraction introduces at most
a simple pole at the other puncture. A polynomial leading term
of degree $j\geq M$ is removed using $t^{j+1}$, whose leading
coefficient is $j+1+a+b$. These reductions terminate.

Conversely,
if $Df\in W_M$, an excessive pole or polynomial degree of $f$
would give an excessive order in $Df$ with one of these nonzero
coefficients. Thus $f\in V_M$, which also proves injectivity
on cohomology.

A nonzero rational horizontal function has orders $-a$ and
$-b$ at zero and one, so both residues must be integral. When
they are, $t^{-a}(t-1)^{-b}$ is a horizontal function and spans
the kernel. The ratio of any two horizontal functions is constant.
The dimension difference of the finite complex therefore gives
shape dimensions $(1,2)$ when $a,b\in\mathbb Z$, and $(0,1)$
otherwise. Tensoring with the two retained scale classes proves
the table. For integral $\lambda$, integrality of $a_{12}$ and
$a_{23}$ is equivalent to integrality of all three residues.
\end{proof}


Thus pairwise global de Rham cohomology cannot replace the full
coefficient diagram: it would discard this ternary class.
The example concerns global cohomology of the displayed connection.
It does not contradict coefficient-level factorization with its
actual maps.

\section{The center of the puncture category}

Return to the regular module and the \(r\) distinct closed coefficients
of Theorem~\ref{cbc:ac:star}.  Their algebra \(A_r\) has identities
\(e_0,\ldots,e_r\) and degree-one arrows
\(\alpha_i:0\to i\), \(\beta_i:i\to0\).
Every product of two arrows is zero.
Write $O=e_0$ for the regular object and $\mathcal L$ for
the finite twisted complex representing localization away from
all $r$ punctures:
$\mathcal L=(\bigoplus_{i=0}^r e_i,d_{\mathcal L}=\sum_i\beta_i)$,
as in Theorem~\ref{cbc:ac:marked-model}.

Put \(S=\bigoplus_{i=0}^rke_i\).  The normalized \(S\)-relative
two-sided bar of \(A_r\) is a resolution of its diagonal bimodule.
Indeed the unnormalized bar contracts by unit insertion, and
separability of \(S\) permits normalization.  Write
\(C(A_r)=\operatorname{Hom}_{A_r^e}(\operatorname{Bar}_S(A_r),A_r)\)
with the full product totalization.  It carries its cup product
and bar insertion operations.

Let \(Q_r^\vee\) have the same vertices and one reversed dual arrow
for each \(\alpha_i,\beta_i\).  Its path-length completion is
\(\widehat{kQ_r^\vee}\).  Define
\[
 \mathcal P=\prod_{i=0}^r e_i\widehat{kQ_r^\vee}e_i,
 \qquad
 \mathcal V=\operatorname{Der}^{\mathrm{cont}}_S(\widehat{kQ_r^\vee}).
\]
A derivation in \(\mathcal V\) is specified by a formal path series
with the same endpoints as each arrow.  These paths have positive
odd length, so every such assignment is continuous.

\begin{theorem}\label{cbc:rc:mpc:thm:centre}
The full cochain center of \(\operatorname{Perf}(A_r)\) is the
two-term complex
\begin{equation}\label{cbc:rc:mpc:eq:centre-complex}
 \mathcal P\xrightarrow{\operatorname{Inn}}\mathcal V,
 \qquad \operatorname{Inn}(p)(x)=xp-px,
\end{equation}
in degrees zero and one.  For \(r\geq2\), its cohomology is
\[
 \operatorname{HH}^0(A_r)=k,\qquad
 \operatorname{HH}^1(A_r)=\prod_{m\geq0}V_{r,m},\qquad
 \operatorname{HH}^q(A_r)=0\quad(q\geq2),
\]
where \(V_{r,m}\) is the quotient of the block of arrow-output
cochains on words of length \(2m+1\), and
\begin{equation}\label{cbc:rc:mpc:eq:dimensions}
 \dim V_{r,0}=r,\qquad
 \dim V_{r,m}=2(r-1)r^m\quad(m\geq1).
\end{equation}
The cup product of two degree-one classes is zero.  With the
negative-transpose identification of arrow-output cochains with
suspended derivations, the Gerstenhaber bracket is the commutator
of derivations modulo inner derivations.
\end{theorem}
\begin{proof}
Every radical input has degree one, hence every suspended input
has degree zero.  Scalar outputs occur on closed even words and
have total degree zero.  Arrow outputs occur on odd words with
matching endpoints and have total degree one.  There are no other
outputs.

The radical products vanish, so every interior face of
the Hochschild differential vanishes.  Its endpoint faces are
precisely \(xp-px\), with the totalization convention of
\eqref{cbc:rc:eq:cup}.  This proves \eqref{cbc:rc:mpc:eq:centre-complex}.
Cup products concatenate words and multiply outputs. Two arrow
outputs multiply to zero.  Inserting cochains into cochains gives
the stated derivation commutator with the chosen transpose sign.

At positive even input length \(2m\), the closed-path space has
dimension \(2r^m\). There are \(r^m\) paths from the central vertex and
\(r^m\) in total from the leaves.  The next odd block has dimension
\(2r^{m+1}\).  At length zero the dimensions are \(r+1\) and \(2r\).

The block differential is explicit.  For an odd path, deleting the first edge gives the closed-path column
with coefficient \(+1\). Deleting the last gives the column with
coefficient \(-1\).  A nonclosed remainder contributes zero.

A kernel element is central in the completed path algebra.  Its
central-vertex component commutes with the \(r\) free return loops.
For \(r\geq2\), a homogeneous word polynomial commuting with two
distinct free generators is scalar. Commutation with the first
restricts it to powers of that generator. Commutation with
the second removes every positive power.

Apply this at each word
length.  The arrow equations force the same scalar at every leaf.
Thus the kernel has dimension one at length zero and vanishes at
every positive length.  Subtracting the ranks gives
\eqref{cbc:rc:mpc:eq:dimensions}.  The differential acts independently on
each finite block, so its image is the product of their images
and is closed.  Its cokernel is exactly the displayed product.
\end{proof}

Order the paths and perform finite elimination in each block of
\eqref{cbc:rc:mpc:eq:centre-complex}. Choose pivot columns for the image
and extend them to a basis of the target. Inverting the differential
on its pivot complement constructs a continuous contraction.
Split off the common identity in length zero.  These choices give
\(dh+hd=1-i\pi\) with the usual side conditions, coefficientwise
in path length.

Choosing the quotient representatives accordingly
gives an associative dg model
\[
 k\oplus\left(\prod_{m\geq0}V_{r,m}\right)[-1]
\]
with zero differential and square-zero positive ideal.  Its full
higher center operations remain on the bimodule bar.  For \(r=1\),
the two cyclic placements of the unique return loop instead give
the logarithmic center \(k[[z]]\otimes\Lambda(\eta)\) of
Theorem~\ref{cbc:rc:thm:centre}, with \([\eta,z]=-z\).

Restriction to the full subcategory on \(e_0,e_i\) gives an actual
map \(C(A_r)\to C(A_1)\).  For \(r\geq2\), its degree-zero map
is the inclusion \(k\to k[[z]]\), and its degree-one map is onto.
To see the latter, extend each arrow-output cochain on those two
vertices by zero on all other words.  All degree-one cochains are closed,
so it remains a closed lift before taking the quotient.

The marked objects also receive this action.  Extend a cochain to
finite twisted complexes by inserting their twisting differentials
between the matrix entries of its inputs, including the two ends.
For \(\mathcal L\), the twisting arrows go from a leaf to the central
vertex, so two consecutive insertions cannot compose.  There is
at most one insertion in each gap.  The extension is therefore
finite at every arity and defines the action on \(O,\mathcal L\).

The bar identities give its differential and higher composition
identities.  This construction uses the specified full dg
subcategory and its twisted complexes, including the ternary
operation in \eqref{cbc:ac:ternary}.

\section{A scale and two punctures}
\label{cbc:rc:scale-shape}

On the chart \(z_1-z_3=u\ne0\), Theorem~\ref{cbc:ac:primitive-chart}
identifies the category generated by the primitive coefficient with
\(\operatorname{Perf}(E\otimes_k A_2)\), where
\(E=\Lambda(\rho)\), \(|\rho|=1\), and \(d\rho=0\).
Here \(\rho\) is represented by \(du/u\).
The shape coordinate \(t=(z_2-z_3)/u\) retains its collisions at
\(0\) and \(1\).  We now compute this category's center and an actual
deformation in its degree-two direction.

Put \(V_2=\prod_{m\geq0}V_{2,m}\). Its block dimensions are
\(2,4,8,16,\ldots\).  The scale Hochschild complex is
\(E[[x]]\), with \(|x|=0\).  Tensor the two bimodule bars at finite
multiweight and then take the product over both weights.  Their
blockwise contractions give the center cohomology on this chart:
\begin{equation}\label{cbc:rc:mpc:eq:chart-centre}
\begin{aligned}
 H^0&=k[[x]],\\
 H^1&=k[[x]]\rho\ \oplus\ k[[x]]\widehat\otimes V_2,\\
 H^2&=k[[x]]\rho\widehat\otimes V_2.
\end{aligned}
\end{equation}
Here the completed tensor product is the product of the finite
multiweight tensor factors.

The tensor bars retain the products
and higher center operations.  The chart includes the two shape
collisions \(t=0,1\).  On its overlap with the second scale chart,
the coordinates are \(t'=1/t\), \(u'=ut\).  Extending the category
through \(u=0\) requires the remaining attachment and comparison
maps on that chart.

\subsection{A product deformation}

Define even derivations \(D_u\) of \(E\) and \(D_t\) of \(A_2\)
by
\[
 D_u(\rho)=\rho,\qquad D_t(\alpha_1)=\alpha_1,
\]
with \(D_t\) zero on the other arrows.  Both kill the identities.
Extend them to their separate factors of \(E\otimes A_2\), whose
product \(\mu\) uses the Koszul sign.
Their commuting additive weights determine the associativity equation.

\begin{proposition}\label{cbc:rc:mpc:prop:twist}
Over \(k[[\varepsilon]]\), the product
\[
 \mu_\varepsilon=\mu\exp(\varepsilon D_u\otimes D_t)
\]
is associative and unital, with zero differential and the same
idempotents.  Its first-order Hochschild class is nonzero in
\(\operatorname{HH}^2(E\otimes A_2)\).
\end{proposition}
\begin{proof}
Every formal coefficient is a defined bilinear map.  On weight
eigenvectors the product acquires the factor
\(\exp(\varepsilon w_u(a)w_t(b))\).  Additivity of the two weights
gives
\[
 w_u(a)w_t(b)+(w_u(a)+w_u(b))w_t(c)
 =w_u(b)w_t(c)+w_u(a)(w_t(b)+w_t(c)),
\]
which proves associativity, including the original Koszul signs.
The identities have weight zero.

The first-order cocycle is the
cup product of the two derivation classes.  The scale class is
\(\rho x\ne0\) in \(\operatorname{HH}^1(E)\).  The shape class is
nonzero in \(V_{2,0}\): an inner diagonal derivation scales the
two directions of an edge oppositely, whereas \(D_t\) scales only
\(\alpha_1\).  The finite-multiweight K\"unneth calculation used in
\eqref{cbc:rc:mpc:eq:chart-centre} makes their tensor cup nonzero.
\end{proof}

For the odd elements \(X=\rho\otimes1\) and
\(Y=1\otimes\alpha_1\), the change is visible as
\[
 X*_\varepsilon Y+Y*_\varepsilon X
 =(e^\varepsilon-1)\rho\otimes\alpha_1.
\]
A realization by native chiral fields must in addition preserve
their collision products, vacuum maps and higher substitutions.
The categorical center computed here represents the bimodule
endomorphism problem on its stated coefficient category. Its
comparison with the full pointed transverse-action problem
requires those additional maps.
\section{Transport of Hochschild operations}

The path substitutions of Theorem~\ref{cbc:ac:braid-substitution}
act on the same completed reverse path algebra \(\widehat P\).
Retain \(\mathcal P=\prod_{i=0}^r e_i\widehat P e_i\), and let
\(\operatorname{Der}^{\mathrm{cont}}_S(\widehat P)\) denote
continuous derivations fixing \(S\).  Full normalized
Hochschild cochains have only total degrees zero and one:
closed even paths have identity outputs, and odd paths have
arrow outputs.  The only differential terms are the two
endpoint actions, since products of radical arrows vanish.
Their model is
\[
 C^\bullet(A_r,A_r)=
 [\mathcal P\longrightarrow
   \operatorname{Der}^{\mathrm{cont}}_S(\widehat P)],
 \qquad f\longmapsto(a\longmapsto af-fa),
\]
with product totalization in path length.

For its Lie structure one shifts Hochschild cochains by one.
Changing the sign on the resulting degree \(-1\) term gives
the convention
\[
 \mathfrak g_r^{-1}=\mathcal P,
 \qquad\mathfrak g_r^0=
          \operatorname{Der}^{\mathrm{cont}}_S(\widehat P),
 \qquad d f=-\operatorname{ad}_f.
\]
The bracket of two derivations is their commutator,
\([\delta,f]=\delta(f)\), and the bracket of two
degree \(-1\) elements is zero.

\begin{proposition}\label{cbc:rc:centre-transport}
Let \(\Phi=\sigma_i^2\) be a pure substitution and
\(\mathcal L=\log\Phi\), as in Lemma~\ref{cbc:ac:log-derivation}.
Then \(\Phi\) acts on this dg Lie complex by
\[
 f\longmapsto\Phi(f),\qquad
 \delta\longmapsto\Phi\delta\Phi^{-1}.
\]
Its infinitesimal action is \(\mathcal L\) on \(\mathcal P\)
and \([\mathcal L,-]\) on derivations.  The corresponding
logarithmic connection is defined on the same completion.
\end{proposition}

\begin{proof}
Conjugation takes \(\operatorname{ad}_f\) to
\(\operatorname{ad}_{\Phi(f)}\), preserves commutators, and
intertwines the evaluation action on \(\mathcal P\).
It therefore preserves the differential and all brackets.
Differentiating \(\Phi^u\) gives the infinitesimal formulas.
A relative derivation sends an arrow to positive odd lengths,
whereas \(\mathcal L\) raises length by at least two.
Its commutator therefore raises the length filtration on
derivations, proving continuity and convergence.
\end{proof}

The full unit-retaining bar carries the Hochschild cup and
brace operations.  Its transport is obtained from the actual
strictly unital higher automorphism and its associated bimodule
bar.  The two-term model above specifies the dg Lie action.
Strict cup invariance of an arbitrary smaller presentation is
an additional condition.


This transport acts on the coefficient category and its Hochschild
operations.  Its identification with a field-theoretic transport would
require maps from the ordered collision coefficients which preserve
both vacuum insertions and every collision substitution.  The relative
collision calculation below will supply a normalization that such maps
must preserve.

\section{An analytic coefficient dual for the regular Heisenberg bar}
\label{cbc:rc:an:sec:analytic-dual}

The two-sided bar above computes diagonal-bimodule endomorphisms.
We next use an augmented bar to construct a spectral coefficient dual.
A nonzero Heisenberg charge requires a coefficient space on which
translation and the collision connection still act. Completing only
the charge coalgebra does not address the positive modes. We construct
a coefficient dual which retains these modes and admits every complex
charge. At fixed positive energy, the bar homotopy uses finitely many mergers,
with a bound independent of the charge degree.

Throughout this section the ground field is \(\mathbb C\), the level
\(\kappa\) is nonzero, and degrees are cohomological. Consider the
augmented commutative dg algebra
\begin{equation}\label{cbc:rc:an:eq:point}
 A_\kappa=\mathbb C[x_n\mid n\geq0]\otimes
             \Lambda(\xi_n\mid n\geq0),\qquad
 |x_n|=0,\quad |\xi_n|=1,\quad
 dx_n=\kappa(n+1)\xi_{n+1},\quad d\xi_n=0.
\end{equation}
Its augmentation kills the generators. The multiplication is the
displayed graded-commutative polynomial multiplication.
Assign energy \(n+1\) to
\(x_n\), and energy \(n\) to \(\xi_n\). The differential preserves
energy, and \(\xi_0\) is the only generator of energy zero.

\subsection{Growth in the charge direction}

Let \(\Gamma(f_0)\) have divided-power basis \(\gamma_n(f_0)\),
\(n\geq0\), with
\(\Delta\gamma_n(f_0)=\sum_{i+j=n}\gamma_i(f_0)\otimes
\gamma_j(f_0)\). For an integer \(R\geq1\), set
\[
 G_R=\left\{a=(a_n)_{n\geq0}:
       |a|_R:=\sup_n |a_n|R^{-n}<\infty\right\},
 \qquad G=\varinjlim_R G_R.
\]
The limit has the locally convex inductive-limit topology. A sequence
represents \(\sum_n a_n\gamma_n(f_0)\). We call \(G\) the analytic moment space. For several
factors we use arrays bounded by \(R^{n_1+\cdots+n_r}\) for some
\(R\), with the same inductive-limit convention.

\begin{lemma}\label{cbc:rc:an:lem:entire}
The continuous linear dual of \(G\) is the space
\(\mathcal E=\mathcal O(\mathbb C)\) of entire functions. The pairing
and its defining seminorms are
\begin{equation}\label{cbc:rc:an:eq:pairing}
 \left\langle\sum_n b_n t^n,a\right\rangle
       =\sum_n b_na_n,
 \qquad \|b\|_R=\sum_n |b_n|R^n\quad(R\geq1).
\end{equation}
Under this pairing, the coproduct \((\Delta a)_{ij}=a_{i+j}\)
is dual to multiplication of entire functions. The sequence
\(a_n=q^n\), for any \(q\in\mathbb C\), represents evaluation at
\(q\).
\end{lemma}
\begin{proof}
The series in \eqref{cbc:rc:an:eq:pairing} is absolutely convergent on
\(G_R\), with absolute value at most \(\|b\|_R|a|_R\).
Conversely, let \(\lambda\) be continuous on \(G\), and let
\(b_n=\lambda(\gamma_n(f_0))\). Continuity on \(G_R\), applied to
finite sequences with independently chosen phases and absolute values
\(R^n\), gives \(\sum_n|b_n|R^n<\infty\). For \(a\in G_R\), its
finite truncations converge to \(a\) in \(G_S\) whenever \(S>R\).

Consequently \(\lambda(a)=\sum_n b_na_n\). No additional functional
on the tails occurs. The coproduct assertion is the Cauchy product,
and \(|q^n|\leq R^n\) for \(R\geq\max(1,|q|)\) proves the last
assertion.
\end{proof}

Let \(\overline A_\kappa\) be the augmentation ideal. In the normalized
bar \(BA_\kappa=\bigoplus_{p\geq0}(s\overline A_\kappa)^{\otimes p}\),
suspension has degree \(-1\): \(|[a]|=|a|-1\).
We use the differential \(D=d_0+b\), where
\[
 d_0[a]=- [da],\qquad b[a|a']=(-1)^{|a|}[aa'].
\]
Extension to a word uses the Koszul sign of the suspended prefix.
The coproduct is deconcatenation, including the two empty prefixes.

Fix the usual monomial basis of \(A_\kappa\), with the odd factors
in increasing index order. Write \(\mathcal W_{E,n}\) for the bar
words of energy \(E\) in which \(\xi_0\) occurs a total of \(n\)
times. Both differentials preserve \(E\) and \(n\).

\begin{definition}\label{cbc:rc:an:def:bar-growth}
For each \(E\geq0\), let
\[
 B^{\mathrm{an}}_E
 =\varinjlim_{R\geq1}
 \left\{(a_w)_{w\in\coprod_n\mathcal W_{E,n}}:
       \sup_{n,w\in\mathcal W_{E,n}}|a_w|R^{-n}<\infty\right\},
 \qquad B^{\mathrm{an}}=\prod_{E\geq0}B^{\mathrm{an}}_E.
\]
Each step carries the displayed supremum norm. The product is taken
in locally convex spaces, degree by degree. For \(r\) coefficient
factors, the corresponding target is the product over
\((E_1,\ldots,E_r)\) of arrays indexed by
\(\mathcal W_{E_1,n_1}\times\cdots\times\mathcal W_{E_r,n_r}\)
and bounded by \(R^{n_1+\cdots+n_r}\) for some \(R\).
\end{definition}

At energy zero this definition gives \(B^{\mathrm{an}}_0=G\), since
the only word is \([\xi_0|\cdots|\xi_0]\). At positive energy the
number of nonzero-energy inputs is at most \(E\). If a word has
length \(p\), then
\begin{equation}\label{cbc:rc:an:eq:length}
 n\leq p\leq n+E,\qquad
 |\mathcal W_{E,n}|\leq C_E(n+E+1)^E
\end{equation}
for a constant \(C_E\). Indeed, each positive-energy input can contain
\(\xi_0\) at most once. There are finitely many possible positive
monomials at energy at most \(E\), and at most \(E\) places among
\(n+E\) positions to insert them. These observations also bound the
cohomological degrees at each fixed \(E\).

\subsection{The completed bar comparison}

Define the dg coalgebra
\[
 C_\kappa=\Lambda(e_n\mid n\geq0)\otimes
                 \Gamma(f_n\mid n\geq0),\qquad
 |e_n|=-1,\quad |f_n|=0,
 \quad de_n=-\kappa(n+1)f_{n+1},\quad df_n=0.
\]
The generators are primitive. The differential is a derivation for
the divided-power product. Give \(e_n,f_n\) energies \(n+1,n\),
respectively.

Let \(C^+_E\) be the energy-\(E\) part generated by
\(e_n\) and \(f_n\), \(n\geq1\) for the latter. It is finite
dimensional. Put
\[
 C^{\mathrm{an}}_\kappa
       =\prod_{E\geq0}(G\otimes C^+_E).
\]
Use the moment-array coefficient powers defined above for its
coproduct. Finite dimensionality makes the tensor topology in each
factor unambiguous.

On a fixed-energy dual we use the seminorms
\(\sum_{n,w}|b_w|R^n\), where \(w\) ranges over the chosen basis
at charge degree \(n\). These seminorms specify its Fr\'echet topology.

\begin{theorem}\label{cbc:rc:an:thm:bar}
The differentials and coproducts extend continuously to
\(B^{\mathrm{an}}\) and \(C^{\mathrm{an}}_\kappa\). The signed shuffle
map extends to a continuous dg coalgebra map
\begin{equation}\label{cbc:rc:an:eq:shuffle}
 \iota:C^{\mathrm{an}}_\kappa\longrightarrow B^{\mathrm{an}},
 \qquad e_n\longmapsto[x_n],\quad f_n\longmapsto[\xi_n].
\end{equation}
It is a quasi-isomorphism. Both complexes admit continuous contractions
onto their common energy-zero subcomplex \(G\).

Topologize their continuous duals by the locally convex direct sum of
the fixed-energy duals. Then \(\iota\) induces a continuous
quasi-isomorphism of dg algebras
\begin{equation}\label{cbc:rc:an:eq:dual-map}
 D^{\mathrm{an}}_\kappa:=(B^{\mathrm{an}})'_{\mathrm{cont}}
 \longrightarrow
 Y^{\mathrm{an}}_\kappa
 :=\mathcal O(\mathbb C_{B_0})\otimes
      \mathbb C[B_n\mid n\geq1]\otimes
      \Lambda(C_n\mid n\geq0),
\end{equation}
where \(|B_n|=0\), \(|C_n|=1\), and
\begin{equation}\label{cbc:rc:an:eq:dual-d}
 dB_n=\kappa nC_{n-1}\quad(n\geq1),
 \qquad dB_0=dC_n=0.
\end{equation}
These are the generators of Section~\ref{cbc:cf:line}, with \(d=Q\).
The energy of \(B_n\) is \(n\), and that of \(C_n\) is
\(n+1\). Multiplication is continuous on each pair of finite-energy
stages.
\end{theorem}
\begin{proof}
The matrices of \(d_0\) and \(b\) on \(\mathcal W_{E,n}\) have
polynomial growth in \(n\). Their entries are bounded in terms of
\(E\) and \(\kappa\). The number of entries is polynomially bounded
by \eqref{cbc:rc:an:eq:length}. A polynomial times \(R^n\) is bounded by a
constant times \(S^n\) for every \(S>R\). Thus both operators extend
continuously at fixed energy, and hence to the product.

A coordinate
of the deconcatenation coproduct is the coordinate of the concatenated
word. The moment growth bound is unchanged, proving continuity into
the stated coefficient powers. The same argument applies to
\(C^{\mathrm{an}}_\kappa\).

We construct the bar contraction directly. Let \(\pi:A_\kappa\to
\Lambda(\xi_0)\) kill the positive-energy generators, and let \(i\)
be the inclusion. On a monomial, let \(N\) count all factors among
\(x_n\) and \(\xi_{n+1}\), with multiplicity. The derivation
\[
 h_A=\sum_{n\geq0}\frac{x_n}{\kappa(n+1)}
                   \frac{\partial}{\partial\xi_{n+1}}
\]
has degree \(-1\) and satisfies \(dh_A+h_Ad=N\). Define
\(H_A=h_A/N\) on \(N>0\), and zero on \(\Lambda(\xi_0)\).
Then \(dH_A+H_Ad=1-i\pi\), \(H_A^2=0\), and \(H_A\) preserves
energy and the number of occurrences of \(\xi_0\).

On the tensor bar with differential \(d_0\), the tensor contraction
\(H\) applies \(-sH_As^{-1}\) to the first positive-energy input.
It is zero if there is no such input. All preceding inputs are
\([\xi_0]\), of degree zero, so this description includes the prefix
sign. If \(I,P\) are the induced inclusion and projection of
\(B\Lambda(\xi_0)\), then
\[
 d_0H+Hd_0=1-IP,\quad H^2=0,\quad HI=PH=0,
 \qquad bI=Pb=0.
\]

The operator \(bH\) lowers bar length by one and preserves \(E,n\).
Equation~\eqref{cbc:rc:an:eq:length} therefore gives
\((bH)^{E+1}=0\) on the entire energy-\(E\) part, independently
of \(n\). Consequently the finite formula
\begin{equation}\label{cbc:rc:an:eq:contraction}
 H'_E=H\sum_{j=0}^{E}(-bH)^j
\end{equation}
satisfies \(DH'_E+H'_ED=1-IP\).

To check the identity at \(E>0\),
put \(U=1+bH\). The relations above give
\(d_0U=Ud_0-b\). Hence
\[
 DHU^{-1}+HU^{-1}D
 =1+HU^{-1}b(1-U^{-1})=1,
\]
since \(1-U^{-1}=bHU^{-1}\) and \(b^2=0\).
At \(E=0\), both sides of the contraction identity vanish.

For fixed \(E\), the matrices in \eqref{cbc:rc:an:eq:contraction} are
bounded by \(K_{E,\kappa}(n+E+1)^{s_E}\), for some integer \(s_E\).
This follows from the finite sum, the bounded entries of \(H_A\),
and \eqref{cbc:rc:an:eq:length}. They therefore carry every growth step
\(R\) continuously into every larger step \(S\). The contractions
extend to \(B^{\mathrm{an}}\).

On \(C^{\mathrm{an}}_\kappa\), the
contracting pairs \((e_n,f_{n+1})\) give a contraction on each finite
space \(C^+_E\). At positive energy it is obtained by dividing the
usual Koszul homotopy by the number of positive generators. Tensoring
with \(G\) and taking the product preserves these identities.

In \eqref{cbc:rc:an:eq:shuffle}, \(\gamma_n(f_0)\) maps to the word of
\(n\) copies of \(\xi_0\), with no factorial. More generally, a
basis monomial maps to the signed sum of its distinct shuffles.
Exchanging two \(e\)'s reverses the sign. The adjacent multiplication
terms cancel by graded commutativity, and the internal differential
has the sign \(de_n=-\kappa(n+1)f_{n+1}\). Thus this is a dg
coalgebra map.

At fixed positive energy the shuffle coefficients have
polynomial growth in \(n\), so the map extends continuously. On energy
zero it is the identity, and both positive-energy complements contract.
This proves the quasi-isomorphism.

A continuous functional on a product of locally convex spaces depends
on finitely many coordinates: a neighborhood controlling it leaves
all remaining coordinates unrestricted. At each energy the proof of
Lemma~\ref{cbc:rc:an:lem:entire} identifies the dual with coefficient
families whose weighted absolute sums converge for every \(R\).
The polynomial bound in \eqref{cbc:rc:an:eq:length} does not affect this
condition. Transposing the contractions gives continuous contractions
of these Fr\'echet spaces, and then of their direct sums. In particular,
transposition preserves the quasi-isomorphism here by an explicit
homotopy, without an exactness assumption on continuous duality.

For \(C^{\mathrm{an}}_\kappa\), its dual is exactly the right side
of \eqref{cbc:rc:an:eq:dual-map}. The convolution product is polynomial in
the positive modes and the ordinary entire-function product in
\(B_0\). The dual differential gives \eqref{cbc:rc:an:eq:dual-d}.
The estimate \(\|fg\|_R\leq\|f\|_R\|g\|_R\), together with the
finite number of positive monomials at fixed energy, proves the last
continuity assertion. The identical weighted-sum estimate proves it
for the bar dual.
\end{proof}

The bar contraction satisfies a sharper continuity statement.
At fixed energy it preserves each charge-growth radius separately.
Write $\mathcal B_{E,R}$ for the full space of coefficient arrays
on $\coprod_n\mathcal W_{E,n}$ with norm
\[
 \|v\|_{E,R}=\sup_{n,w\in\mathcal W_{E,n}}|v_w|R^{-n},
 \qquad R\geq1.
\]
This is the bounded-array space in
Definition~\ref{cbc:rc:an:def:bar-growth}, not the norm closure of the
finite arrays.

\begin{theorem}\label{cbc:rc:an:thm:fixed-radius}
For $E\geq1$, the operator $H'_E$ of
\eqref{cbc:rc:an:eq:contraction} extends to a bounded degree-minus-one
operator on $\mathcal B_{E,R}$ with
\[
 DH'_E+H'_ED=1,\qquad
 \|H'_E\|\leq\frac1{|\kappa|}
 \sum_{j=0}^{E}\left(\frac{E2^{E+1}}{|\kappa|}\right)^j.
\]
The radius $R$ is unchanged. At energy zero the differential
vanishes on all arrays of words in $[\xi_0]$, including the empty
word. No bound uniform in $E$ is asserted.
\end{theorem}
\begin{proof}
We bound incoming row sums in the bar basis: fix an output word
and sum absolute coefficients over its possible inputs.
For one output monomial $x^a\xi_S$, an incoming internal
differential replaces $x_j$ by $\xi_{j+1}$ and requires
$j+1\in S$. Its coefficient has magnitude
$|\kappa|(j+1)(a_j+1)$.
Their sum is at most
\[
 |\kappa|\left(\sum_j(j+1)a_j+
                    \sum_{s\in S\setminus\{0\}}s\right).
\]
Summing over the word gives $\|d_0\|\leq|\kappa|E$.

For $H$, fix the first positive-energy output input, with
positive factor count $N$. An incoming term replaces one of
its $x_j$ by $\xi_{j+1}$. It requires $a_j>0$ and
$j+1\notin S$. Its coefficient has magnitude
$1/(|\kappa|(j+1)N)$. There are at most $N$ eligible indices,
so $\|H\|\leq1/|\kappa|$.

For $b$, an incoming term splits one output monomial into an
ordered pair of nonconstant monomials. The output must have
positive energy, because $\xi_0$ has no such nonzero
factorization. There are at most $E$ eligible output inputs.
If one has polynomial exponents $a_j$, $r$ positive-energy odd
factors, and possibly $\xi_0$, the number of ordered
factorizations is at most
\[
 \prod_j(a_j+1)\,2^{r+\mathbf1_{0\in S}}
 \leq2^{\sum_j a_j+r+\mathbf1_{0\in S}}
 \leq2^{E+1}.
\]
This count includes forbidden constant factors and therefore
remains an upper bound after normalization. Each contribution
has magnitude one, giving $\|b\|\leq E2^{E+1}$.

All these matrices preserve $n$ and have finite rows. The
weights $R^{-n}$ therefore cancel, so the bounds hold on the
full bounded-array space without a density argument.
The algebraic contraction identities extend coefficientwise.
Since $bH$ lowers length and $n\leq p\leq n+E$,
$(bH)^{E+1}=0$ on this entire space. The finite formula
\eqref{cbc:rc:an:eq:contraction} gives the asserted identity and norm
bound by submultiplicativity.
\end{proof}

For multilinear operations on these duals, use the bornology defined
as follows. A bounded set lies in finitely many energy degrees and is
bounded for every entire-function seminorm there. Multiplication
is bounded for this bornology.

Define the \(r\)-factor target \((Y^{\mathrm{an}}_\kappa)^{(r)}\)
as a direct sum over nonnegative energy indices. Each summand is
\(\mathcal O(\mathbb C^r)\) tensored with the corresponding finite space
of positive monomials. Its seminorms on the entire
coefficients are the weighted absolute sums with radii
\((R_1,\ldots,R_r)\). This specifies the coefficient tensor powers
used below.

\begin{corollary}\label{cbc:rc:an:cor:cohomology}
One has
\[
 H^*(Y^{\mathrm{an}}_\kappa)=\mathcal O(\mathbb C_{B_0})
 \quad\text{in degree zero}.
\]
The inclusion of the polynomial algebra \(Y_\kappa\) from
Section~\ref{cbc:cf:line} into \(Y^{\mathrm{an}}_\kappa\) is not a
quasi-isomorphism.
\end{corollary}
\begin{proof}
The degree \(-1\) derivation sending
\(C_n\) to \(B_{n+1}/(\kappa(n+1))\) and killing all
\(B_n\) has commutator with \(d\) equal to the number of positive
generators. Divide by this number off the charge subalgebra. This is
a continuous contraction at each energy. On cohomology the inclusion
is \(\mathbb C[B_0]\hookrightarrow\mathcal O(\mathbb C_{B_0})\),
which is not surjective because \(e^{B_0}\) is not polynomial.
\end{proof}

\subsection{Spectral operations and charge transport}

Define derivations \(T,L\) of degrees \(0,-1\) by
\begin{equation}\label{cbc:rc:an:eq:cartan}
 Tg_n=ng_{n-1}\quad(g_n=B_n\text{ or }C_n),\qquad
 LC_n=B_n/\kappa,\qquad LB_n=0,
\end{equation}
with \(Tg_0=0\). These agree with \eqref{cbc:cf:operators}. Both kill entire functions of \(B_0\).
They lower energy by one whenever they are nonzero. They are
continuous and satisfy \(dL+Ld=T\). Thus
\(\tau_z=\exp(zT)\) is a finite sum on every element, for every
\(z\in\mathbb C\).

Define a continuous algebra map
\(\Delta_z:Y^{\mathrm{an}}_\kappa\to
(Y^{\mathrm{an}}_\kappa)^{(2)}\) by
\begin{equation}\label{cbc:rc:an:eq:coproduct}
 \Delta_z(g_n)=\tau_z(g_n)\otimes1+1\otimes g_n,
 \qquad \Delta_z(f(B_0))=f(B_0\otimes1+1\otimes B_0).
\end{equation}
The counit evaluates \(B_0\) at zero and kills the positive modes.
Continuity of the second formula follows from
\(\|f(t_1+t_2)\|_{R_1,R_2}\leq\|f\|_{R_1+R_2}\).

Let
\[
 c_{nm}(z)=(-1)^n\binom{n+m}{n}z^{-n-m-1},
 \qquad
 r(z)=\sum_{n,m\geq0}c_{nm}(z)
       (C_n\otimes B_m-B_n\otimes C_m).
\]
This degree-one tensor belongs to
\((Y^{\mathrm{an}}_\kappa)^{(2)}[[z^{-1}]]\), with finitely many
terms in each coefficient. Translations of this series are formal
Taylor translations.

Interpret identities involving several positions as rational jet
identities. The coefficients \(c_{nm}(z)\) are the Taylor coefficients
in \(u,v\) of \((z+u-v)^{-1}\). On three positions the coefficient ring is
\[
 \mathbb C[z_{12},z_{23},
     z_{12}^{-1},z_{23}^{-1},(z_{12}+z_{23})^{-1}]
       [[u_1,u_2,u_3]].
\]
Expansion in a chosen position domain is made only after this rational
identity. Numerical convergence requires a further positive-mode completion,
such as the one in Theorem~\ref{cbc:ac:exact-domain}.

\begin{proposition}\label{cbc:rc:an:prop:spectral}
The maps \(\Delta_z,\tau_z\), the counit and \(r(z)\) satisfy the
following spectral identities on the specified coefficient powers:
\begin{gather*}
 dr=0,\qquad (T\otimes1)r=\partial_zr,
 \qquad (1\otimes T)r=-\partial_zr,\\
 (\Delta_w\otimes1)\Delta_z=(1\otimes\Delta_z)\Delta_{z+w},\\
 (\Delta_w\otimes1)r(z)+r_{12}(w)
   =(1\otimes\Delta_z)r(z+w)+r_{23}(z).
\end{gather*}
Moreover, the degree-zero tensor
\begin{equation}\label{cbc:rc:an:eq:connection}
 X(z)=(L\otimes1)r(z)
       =\frac1\kappa\sum_{n,m\geq0}c_{nm}(z)
                                 B_n\otimes B_m
\end{equation}
satisfies \(dX=\partial_zr\). With exterior forms written on the
left, \(r+dz\,X\) satisfies the Maurer--Cartan equation for
\(d+d_z\). On three distinct positions, so does
\(\sum_{i<j}(r_{ij}(z_{ij})+dz_{ij}X_{ij}(z_{ij}))\), where
\(z_{ij}=z_i-z_j\).
\end{proposition}
\begin{proof}
The derivations in \eqref{cbc:rc:an:eq:cartan} are checked on generators.
Their extensions to entire coefficients follow from their vanishing
on \(B_0\). The same argument proves that \(\Delta_z\) is a
cochain map. Both coproduct iterates send \(g_n\) to
\[
 \tau_{z+w}(g_n)\otimes1\otimes1
 +1\otimes\tau_z(g_n)\otimes1
 +1\otimes1\otimes g_n,
\]
and send \(f(B_0)\) to \(f(B_0^{(1)}+B_0^{(2)}+
B_0^{(3)})\). This proves coassociativity on the whole carrier.

The coefficient identities
\[
 (n+1)c_{n+1,m}=\partial_zc_{nm},\qquad
 (m+1)c_{n,m+1}=-\partial_zc_{nm}
\]
prove the translation formulas. The coefficient of
\(C_n\otimes C_m\) in \(dr\) is
\(-\kappa((n+1)c_{n+1,m}+(m+1)c_{n,m+1})=0\).
Primitivity gives
\[
 (\Delta_w\otimes1)r(z)=r_{13}(z+w)+r_{23}(z),\qquad
 (1\otimes\Delta_z)r(z+w)=r_{12}(w)+r_{13}(z+w),
\]
in the rational jet domains. The remaining counit and graded-flip
identities follow on generators from \eqref{cbc:rc:an:eq:coproduct} and
\(c_{nm}(-z)=-c_{mn}(z)\). These arguments give the same identities
on entire coefficients by the displayed seminorm estimate.

Finally \([d,L]=T\) and \(dr=0\) give
\(dX=(T\otimes1)r=\partial_zr\). In the tensor product with
exterior forms, \(d(dz\,X)=-dz\,dX\). Thus
\((d+d_z)(r+dz\,X)=0\). The coefficient algebra is graded
commutative, so the square of this odd element vanishes. The same
calculation applies to the sum on three positions: each differential
vanishes and all graded commutators vanish. Every calculation is
coefficientwise in the specified formal domains.
\end{proof}

With the sign convention of \eqref{cbc:cf:r}, the charge residue is
positive. The positive-mode terms of this superconnection can be removed over
the formal punctured neighborhood of infinity. The remaining charge
connection is retained by the following explicit comparison.

\begin{proposition}\label{cbc:rc:an:prop:gauge}
Put \(\mathscr Y=(Y^{\mathrm{an}}_\kappa)^{(2)}((z^{-1}))\), meaning
Laurent series in \(z^{-1}\) with coefficients in the indicated
coefficient space and a finite upper bound on powers of \(z\).
Each fixed-bound space has the product topology. Their union has the
inductive-limit topology. Define
\begin{equation}\label{cbc:rc:an:eq:gauge}
 b(z)=-\frac1\kappa\sum_{\substack{n,m\geq0\\n+m>0}}
       \frac{(-1)^n}{n+m}\binom{n+m}{n}
       B_n\otimes B_m\,z^{-n-m},\qquad g(z)=e^{b(z)}.
\end{equation}
These are well-defined formal series, with \(g=1+O(z^{-1})\).

Writing \(D=d+d_z\), multiplication by \(g\) is an isomorphism
of complexes from
\[
 \bigl(\mathscr Y\otimes\Lambda(dz),\ D+r+dz\,X\bigr)
 \quad\text{to}\quad
 \left(\mathscr Y\otimes\Lambda(dz),\
 D+\frac{B_0\otimes B_0}{\kappa}\frac{dz}{z}\right).
\]
It is continuous in the coefficientwise formal topology. Projection
\(p\) to energy zero is a quasi-isomorphism from the first complex to
the following complex, where \(t_i=B_0^{(i)}\) are the charge
coordinates:
\begin{equation}\label{cbc:rc:an:eq:charge-complex}
 \left(\mathcal O(\mathbb C^2)((z^{-1}))\otimes\Lambda(dz),\
 d_z+\frac{t_1t_2}{\kappa}\frac{dz}{z}\right).
\end{equation}
A homotopy inverse is \(f\mapsto e^{-b(z)}f\).
\end{proposition}
\begin{proof}
Write \(r=\sum_{j\geq1}r_jz^{-j}\),
\(X=\sum_{j\geq1}X_jz^{-j}\), and
\(b=\sum_{j\geq1}b_jz^{-j}\). Proposition~\ref{cbc:rc:an:prop:spectral}
gives \(dX_{j+1}=-j r_j\). Formula~\eqref{cbc:rc:an:eq:gauge} says
\(b_j=-X_{j+1}/j\), and therefore
\[
 db=r,\qquad
 \partial_z b=X-X_1z^{-1},\qquad
 X_1=\kappa^{-1}B_0\otimes B_0.
\]
The exponential and its inverse exist because \(b\) has strictly
positive \(z^{-1}\)-order. Each coefficient uses finitely many
products. Those products are continuous on fixed-energy stages.

In fact each coefficient of \(g\) is a polynomial in finitely many
positive modes and the two charge variables. Multiplication by it
is continuous on the whole locally convex direct sum. This proves
coefficientwise continuity of multiplication by \(g\) and \(g^{-1}\).
Graded commutativity now gives
\[
 g^{-1}(D+dz\,X_1z^{-1})g
       =D+dz\,X_1z^{-1}+Db=D+r+dz\,X.
\]

Sum the two Koszul homotopies of
Corollary~\ref{cbc:rc:an:cor:cohomology} and divide by the total number of
positive generators away from energy zero. This contraction commutes
with multiplication by \(B_0\otimes B_0\). Extend it to forms
with the Koszul sign and to Laurent coefficients term by term. It
then contracts the second complex onto \eqref{cbc:rc:an:eq:charge-complex}.
Conjugating this contraction by \(g\) gives the asserted contraction
of the first complex. Since \(p(b)=0\), one has \(pg=p\), and its
inclusion is precisely multiplication by \(e^{-b}\) on energy zero.
\end{proof}

\begin{example}[Nonzero charges]\label{cbc:rc:an:ex:charges}
For every \(q\in\mathbb C\), the formula
\[
 \chi_q(f(B_0))=f(q),\qquad
 \chi_q(B_n)=0\ (n\geq1),\qquad
 \chi_q(C_n)=0\ (n\geq0)
\]
defines a continuous dg algebra character. On a pair of charges
\(q,q'\), the superconnection of
Proposition~\ref{cbc:rc:an:prop:spectral} becomes
\[
 \nabla_{q,q'}=d_z+\frac{qq'}{\kappa}\frac{dz}{z}.
\]
Its horizontal sections on a logarithm branch are scalar multiples
of \(z^{-qq'/\kappa}\). Counterclockwise full transport has scalar
\(e^{-2\pi iqq'/\kappa}\), and the half-turn along
\(z(t)=z_0e^{\pi it}\), with the logarithm continued along this path,
has scalar \(e^{-\pi iqq'/\kappa}\). On three charges the horizontal
factor is \(\prod_{i<j}z_{ij}^{-q_iq_j/\kappa}\), with a consistent
choice of logarithms on the universal cover of the configuration
space.

These characters do not commute with translation when \(q\ne0\):
\(\chi_q(TB_1)=q\), whereas \(T\chi_q(B_1)=0\).
Consequently a fixed nonzero charge cannot replace the full algebra
in the spectral coproduct identities. The positive modes in
\eqref{cbc:rc:an:eq:dual-map} are necessary even though their cohomology
vanishes.
\end{example}

The construction above supplies a continuous coefficient dual of the
specified completion of the regular augmented algebra
\eqref{cbc:rc:an:eq:point}. A comparison with the state-dependent bar of the nonaugmented
Heisenberg transverse monoid must preserve the chosen state and
the two-position collision. It must also preserve all three
partitions of three positions. In particular, the ordinary augmentation
of \(A_\kappa\) is not an augmentation of that transverse monoid.
Theorem~\ref{cbc:rc:an:thm:bar} constructs the coefficient dual on the
regular polynomial carrier. This comparison with the transverse
monoid remains a separate mathematical question.

\section{Multiplication from commuting resolutions}
\label{cbc:rc:tower}

The action in~\eqref{cbc:va:cochain-roof} resolves the boundary by
one quartet. To compose two actions, use two quartets. This gives
strict products with different targets. The maps that merge those
targets are quasi-isomorphisms, and their compatibility supplies the
coherent product.

Work over a field $k$ of characteristic zero, on $X=\mathbb A^1_k$,
with $\kappa\in k^\times$. All state spaces in this section are
algebraic. Retain the full ordered collision atlas and its category
$\operatorname{Coalg}^{\mathrm f}$ from
Section~\ref{cbc:va:pointed}. Its diagonal tensor is $\otimes^!$,
with vacuum unit $J$. The symbol $J$ also denotes the Heisenberg
current when it occurs inside a vertex algebra.

Write $\mathcal Z=\mathcal Z_\kappa$ for the cochain vertex algebra
of Section~\ref{cbc:va:cochain-action}, and write $E_{\mathrm q}$
for its quartet algebra. Thus
\[
 [J_\lambda J]=\kappa\lambda,\qquad dJ=\kappa Tc,\quad dc=0,
 \qquad
 [p_\lambda t]=[\iota_\lambda\eta]=1,
 \quad Qt=\eta,\quad Q\iota=p.
\]
The degrees of $t,p,\eta,\iota$ are $0,0,1,-1$.

For $b\geq1$ put
\[
 W_b=V_\kappa\otimes E_{\mathrm q}^{\otimes b},\qquad
 K_b=J+\sum_{j=1}^b p_j,
 \qquad i_b(J)=K_b.
\]
The differential $Q$ fixes $J$ and acts on each quartet as above.
Replacing $J$ by $K_b$ in the creation-mode basis leaves $b$
contractible quartets. Their Euler contraction, as in
Section~\ref{cbc:va:cochain-action}, proves that each $i_b$ is a
quasi-isomorphism. Let $\mathcal Z^{(b)}$ be the commutant of
$i_b(V_\kappa)$ in $W_b$, namely the common kernel of the
nonnegative modes of $K_b$.

If $\phi:[b']\twoheadrightarrow[b]$ is an order-preserving
surjection, let $\phi_*(j)=\min\phi^{-1}(j)$. Define a merge
$m_\phi:W_b\to W_{b'}$ by fixing $J$ and putting
\begin{equation}\label{cbc:rc:tower-merge}
 p_j\longmapsto\sum_{\phi(l)=j}p_l,\qquad
 \iota_j\longmapsto\sum_{\phi(l)=j}\iota_l,\qquad
 t_j\longmapsto t_{\phi_*(j)},\qquad
 \eta_j\longmapsto\eta_{\phi_*(j)}.
\end{equation}
For $a\geq1$ the shift $r_a:W_b\to W_{a+b}$ sends $J$ to $K_a$
and sends a quartet with index $j$ to the quartet with index $a+j$.
Include $W_a$ by its first $a$ quartets. The candidate product is
\begin{equation}\label{cbc:rc:tower-product}
 \mu_{a,b}:\mathcal Z^{(a)}\otimes W_b\longrightarrow W_{a+b},
 \qquad z\otimes w\longmapsto z_{(-1)}r_a(w).
\end{equation}

\begin{proposition}\label{cbc:rc:tower-strict}
The merges are quasi-isomorphisms of dg vertex algebras, satisfy
$m_\phi i_b=i_{b'}$ and compose contravariantly in $\phi$.
They restrict to quasi-isomorphisms
$\mathcal Z^{(b)}\to\mathcal Z^{(b')}$.
The maps~\eqref{cbc:rc:tower-product} are dg vertex maps, preserve
the commutants, and are strictly associative and natural under
concatenation of merges.

The algebra $\mathcal Z^{(b)}$ is freely generated by
\[
 U_1=J-\kappa Tt_1,\qquad
 t_j-t_1\ (2\leq j\leq b),\qquad p_j,\eta_j,\iota_j\ (1\leq j\leq b).
\]
The map $j_b:\mathcal Z\to\mathcal Z^{(b)}$ with
$j_b(J)=U_1$ and $j_b(c)=-\eta_1$ is a quasi-isomorphism.
The iterated product of $b$ cochain factors, together with its
boundary action, sends their generators to
\begin{equation}\label{cbc:rc:tower-currents}
 J^{(i)}\longmapsto U_i=J+\sum_{l<i}p_l-\kappa Tt_i,
 \qquad c^{(i)}\longmapsto-\eta_i,
 \qquad J_{\mathrm{boundary}}\longmapsto K_b.
\end{equation}
\end{proposition}
\begin{proof}
In~\eqref{cbc:rc:tower-merge}, the sum of the $p_l$ in a fiber
pairs with its selected $t_{\phi_*(j)}$ through exactly one term.
The same holds for $\iota$ and $\eta$. All other brackets are
preserved, as are $Q$ and $T$. The fibers partition $[b']$, so
$m_\phi(K_b)=K_{b'}$.

The least point over the least point is the
least point of a composite fiber. This proves composition.
The identity $m_\phi i_b=i_{b'}$ then proves that $m_\phi$ is a
quasi-isomorphism.

For the commutant, write $x_n=J_{-n}\mathbf1$ for $n\geq1$ and
$\tau_n^{(j)}=T^nt_j/n!$ for $n\geq0$. On the creation-mode
basis the relevant modes act by
\[
 (K_b)_0=\sum_j\partial_{\tau_0^{(j)}},\qquad
 (K_b)_n=\kappa n\partial_{x_n}
                  +\sum_j\partial_{\tau_n^{(j)}}\quad(n\geq1).
\]
Replace $x_n$ by $x_n-\kappa n\tau_n^{(1)}$ and each
$\tau_n^{(j)}$, $j>1$, by $\tau_n^{(j)}-\tau_n^{(1)}$.
The displayed modes become the derivatives in the remaining
$\tau_n^{(1)}$. Their common kernel has precisely the generators
in the statement. The change is invertible over $k[T]$, so there
are no additional relations.

On these generators,
\[
 QU_1=-\kappa T\eta_1,
 \quad Q(t_j-t_1)=\eta_j-\eta_1,
 \quad Q\iota_j=p_j.
\]
Thus its complex of creation states is that of $\mathcal Z$,
with $b-1$ contractible pairs of differences and $b$ pairs
$(\iota_j,p_j)$ adjoined. All states are finite polynomials.
Dividing the Euler contraction by positive pair degree contracts
the added pairs. This proves the assertion about $j_b$.

Since every monotone surjection has $\phi_*(1)=1$, one has
$m_\phi j_b=j_{b'}$. Its restriction to the commutants is
therefore a quasi-isomorphism.

The current $K_a$ has level $\kappa$, is closed, and commutes
with the later quartets. Hence $r_a$ is a dg vertex map.
The images of $\mathcal Z^{(a)}$ and $r_a(W_b)$ commute.
Their normally ordered product consequently defines a map of
the tensor product vertex algebra. Both images of the
commutants commute with $K_{a+b}$.

The identities
$r_ar_b=r_{a+b}$ and
$m_{\phi\sqcup\psi}r_a=r_{a'}m_\psi$ prove associativity
and naturality. Applying the successive shifts to $j_1$ gives
\eqref{cbc:rc:tower-currents}.
\end{proof}

The $b+1$ currents $U_1,\ldots,U_b,K_b$ commute with one another
and each has level $\kappa$. The space spanned by $J,p_j,Tt_j$
has the form with $J$ of norm $\kappa$ and each pair $p_j,Tt_j$
hyperbolic. Thus each additional commuting current uses one
additional quartet.

\begin{lemma}\label{cbc:rc:tower-localization}
Let $\mathcal P$ have the nonempty finite ordered sets $[b]$ as
objects and monotone surjections as morphisms. Let $\mathcal D$
be a symmetric monoidal infinity category. Suppose
$Z:\mathcal P^{\mathrm{op}}\to\mathcal D$ is lax monoidal for
concatenation, and $M:\mathcal P^{\mathrm{op}}\to\mathcal D$
has natural associative maps $Z(a)\otimes M(b)\to M(a+b)$.
If both functors carry every arrow to an equivalence, then
$Z(1)$ is a nonunital associative algebra and $M(1)$ its
nonunital left module. Their operations are the roofs
\[
 Z(1)^{\otimes n}\longrightarrow Z(n)\xleftarrow{\sim}Z(1),
 \qquad Z(1)\otimes M(1)\longrightarrow M(2)\xleftarrow{\sim}M(1).
\]
\end{lemma}
\begin{proof}
Encode the pair of functors by the two-color operad for a
nonunital associative algebra and a left module. Each color carries
an index in $\mathcal P^{\mathrm{op}}$. Substitution
concatenates these indices. Passing to symmetric monoidal
envelopes makes the given data a monoidal functor
\cite{LurieHA}*{Proposition 2.2.4.9}.

The map of envelopes that forgets the indices is a cocartesian
fibration. Its fibers are finite products of
$\mathcal P^{\mathrm{op}}$. Mark every fiber arrow, and every
arrow over a base equivalence. Concatenation preserves these
markings. Localization over the base is therefore computed by
localizing the fibers
\cite{Hinich2016Localization}*{Proposition 2.1.4}.

The nerve of $\mathcal P^{\mathrm{op}}$ is contractible because
$[1]$ is initial. Each localized fiber is consequently a point.
The localized envelope is the envelope of the two-color operad.
The original functor inverts the marked arrows, so it factors
through this monoidal localization
\cite{LurieHA}*{Proposition 4.1.7.4}. This gives the algebra,
module, and all their coherences. The unique surjection
$[n]\twoheadrightarrow[1]$ gives the displayed roofs.
\end{proof}

\begin{theorem}\label{cbc:rc:cochain-monoid}
The coalgebra $C_{\mathcal Z}$ has a unital associative
multiplication for $\otimes^!$, with its vacuum as unit.
Its product of $b$ factors is induced by the roof
\[
 \mathcal Z^{\otimes b}\longrightarrow\mathcal Z^{(b)}
 \xleftarrow{\sim}\mathcal Z^{(1)}\xleftarrow{j_1}\mathcal Z,
\]
whose first map is~\eqref{cbc:rc:tower-currents} and whose
middle map merges all quartets. It has a unital coherent action
on $C_{V_\kappa}$ with underlying map
\eqref{cbc:va:cochain-roof}.
\end{theorem}
\begin{proof}
Apply Lemma~\ref{cbc:va:free-field-tensor} to
$Z(b)=C_{\mathcal Z^{(b)}}$ and $M(b)=C_{W_b}$.
Proposition~\ref{cbc:rc:tower-strict} supplies the hypotheses
of Lemma~\ref{cbc:rc:tower-localization}. This first gives
a nonunital algebra and module.

Write $m:W_1\to W_2$ for the merge. Right multiplication by
the vacuum on $\mathcal Z^{(1)}$ is the inclusion into
$\mathcal Z^{(2)}$, which agrees with $m$ after $j_1$.
Left multiplication is $r_1$. Use polynomial forms on the
interval, standing to the left of fields. The following two
families join $mj_1$ to $r_1j_1$:
\begin{align}
 J&\longmapsto J-\kappa T((1-s)t_1+st_2),
 &c&\longmapsto-(1-s)\eta_1-s\eta_2-ds(t_2-t_1),
 \label{cbc:rc:left-unit-one}\\
 J&\longmapsto J+sp_1+ds\,\iota_1-\kappa Tt_2,
 &c&\longmapsto-\eta_2.
 \label{cbc:rc:left-unit-two}
\end{align}

Their images commute with $K_2$, have the brackets of
$\mathcal Z$, and satisfy $dJ=\kappa Tc$, $dc=0$ for the
total differential. For example,
$d(sp_1+ds\,\iota_1)=ds\,p_1-ds\,p_1=0$.
The first endpoint of the second family is the last endpoint
of the first. Lemma~\ref{cbc:va:free-field-tensor} carries
these polynomial homotopies to the coalgebra category.
Since $C_{j_1}$ is an equivalence, both multiplications by the
vacuum are homotopic to the identity.

On the module the vacuum acts by $r_1:W_1\to W_2$.
The equation $r_1i_1=i_2$ makes this a quasi-isomorphism.
The algebra and module are therefore quasi-unital. The
quasi-unit theorems give their unital structures, with a
contractible space of choices
\cite{LurieHA}*{Theorem 5.4.3.5, Proposition 5.4.3.14}.
Transport them along $C_{j_1}$ and $C_{i_1}$.

For the underlying action, the identity
\[
 \mu_{1,1}(j_1(z)\otimes i_1(v))
   =j_1(z)_{(-1)}i_2(v)=m\bigl(\rho(z\otimes v)\bigr)
\]
identifies its roof with~\eqref{cbc:va:cochain-roof}.
Here $m$ fixes $j_1(z)$, since it changes only $p$ and
$\iota$ in the first quartet.
\end{proof}

The tensor comparison used here is over the ground field, with
derived tensor on the coefficient schemes. Ordinary tensor over
a differential parameter algebra has a different behavior.

\begin{example}\label{cbc:rc:parameter-tensor}
Put $\Omega=k[s,ds]$ and consider regular commutative vertex
algebras
\[
 A_+=\Omega[a_n:n\geq0],\qquad A_-=\Omega[b_n:n\geq0],
 \qquad Da_n=ds\,a_n,
 \quad Db_n=-ds\,b_n.
\]
All $a_n,b_n$ have degree zero, their brackets vanish,
$Ta_n=a_{n+1}$, $Tb_n=b_{n+1}$, and $T\Omega=0$.
Both augmentations to $\Omega$ are quasi-isomorphisms.
On a monomial of positive degree $m$ the coefficient
differential is $\partial_s\pm m$, whose inverse on a
polynomial is the terminating series
\[
 (\partial_s+c)^{-1}f
   =\sum_{n\geq0}\frac{(-1)^n}{c^{n+1}}\partial_s^nf
   \quad(c\ne0).
\]
In $A_+\otimes_\Omega A_-$, however, $a_0b_0$ is closed
and cannot be a boundary.

More precisely, give $B=k[a_i,b_j]$ charge $1$ on the $a_i$
and $-1$ on the $b_j$, and write $B_0$ for charge zero.
The inclusion $\Omega\otimes_k B_0\to A_+\otimes_\Omega A_-$
is a quasi-isomorphism. Indeed the coefficient complex of a
charge-$c$ monomial has differential $\partial_s+c$.
It is acyclic for $c\ne0$ and has the constants as its
degree-zero cohomology for $c=0$.

The surviving algebra has the presentation
\[
 B_0=k[z_{ij}:i,j\geq0]/(z_{ij}z_{kl}-z_{il}z_{kj}),
 \qquad z_{ij}=a_ib_j,
 \qquad Tz_{ij}=z_{i+1,j}+z_{i,j+1}.
\]
Every balanced monomial is a product of the $z_{ij}$.
The displayed relations exchange column partners, so every
monomial reduces to the pairing of its sorted row and column
multisets. Distinct pairs of multisets give distinct monomials
in the independent $a_i,b_j$. This proves the presentation,
including for countably many variables, since every relation
uses only finitely many. With $r$ row indices and $s$ column
indices, the degree-$q$ dimension is
$\binom{r+q-1}{q}\binom{s+q-1}{q}$.

The inclusions $\Omega\to A_\pm$ are quasi-isomorphisms from
a $K$-flat module. Hence the comparison from derived to
ordinary tensor has degree-zero cohomology map $k\to B_0$,
$1\mapsto1$. Neither $A_+$ nor $A_-$ is $K$-flat over
$\Omega$: tensoring the other's augmentation with it would
otherwise make the augmentation $B_0\to k$ an isomorphism.
\end{example}

\section{The endomorphisms of a moving vacuum}
\label{cbc:rc:moving-vacuum}

The cochain action has a universal property at one position.
Its tests there are modules for the boundary current. These
tests retain the zero mode and the connection as the position
moves.

Let $\mathfrak h_\kappa$ have generators $J_n$, $n\in\mathbb Z$,
and a central generator $\mathbf c$, with
\[
 [J_m,J_n]=\kappa m\delta_{m+n,0}\mathbf c.
\]
A smooth module has $\mathbf c=1$, and each of its vectors is
killed by all sufficiently large $J_n$. The vacuum is
$V_\kappa=k[x_1,x_2,\ldots]$, where $J_{-n}$ multiplies by
$x_n$, $J_n=\kappa n\partial_{x_n}$ for $n>0$, and $J_0=0$.
Write $\mathfrak b=\bigoplus_{n\geq0}kJ_n$ and
$\mathfrak a_p=\bigoplus_{n\geq p}kJ_n$.

For a smooth module $M$, its continuous cochains are
\begin{equation}\label{cbc:rc:continuous-cochains}
 C^\bullet(M)=M\otimes\Lambda(\xi_0,\xi_1,\ldots),
 \qquad d(m\omega)=\sum_{n\geq0}J_nm\,\xi_n\wedge\omega.
\end{equation}
The sum is finite on each element. These cochains are the
filtered union of the finite-dimensional Lie algebra
cochains of $\mathfrak b/\mathfrak a_p$ on
$M^{\mathfrak a_p}$.

\begin{lemma}\label{cbc:rc:smooth-ext}
The complex~\eqref{cbc:rc:continuous-cochains} computes
$R\operatorname{Hom}(V_\kappa,M)$ in the category of smooth
modules. In particular
\[
 \operatorname{Ext}^*(V_\kappa,V_\kappa)=\Lambda(\xi_0).
\]
The class $\xi_0$ is represented by the nonsplit extension
$V_\kappa\otimes k[\epsilon]/(\epsilon^2)$ on which $J_0$
acts by $\epsilon$.
\end{lemma}
\begin{proof}
The vacuum is induced from the trivial $\mathfrak b$-module,
so its homomorphism functor is $M\mapsto M^{\mathfrak b}$.
The cochain functors are exact. For
$S_p=k[J_0,\ldots,J_{p-1}]$, induction to smooth Heisenberg
modules is exact by the creation-mode basis and is left
adjoint to $M\mapsto M^{\mathfrak a_p}$.

Thus $Y^{\mathfrak a_p}$ is injective over $S_p$ when $Y$
is injective among smooth modules. Its finite Koszul cochains
are acyclic in positive degrees. Filtered unions are exact,
so the continuous cochains have the same property on $Y$.
They therefore compute the derived invariant functor.

For the vacuum the differential is
$\sum_{n\geq1}\kappa n\xi_n\partial_{x_n}$.
The degree-minus-one derivation
$h(\xi_n)=x_n/(\kappa n)$ for $n\geq1$, with $h(\xi_0)=h(x_n)=0$,
satisfies $dh+hd=N$, where $N$ counts the positive-index
generators. Dividing by $N$ contracts positive degree and
leaves $\Lambda(\xi_0)$. The proposed extension is nonsplit
because $J_0$ vanishes on a split extension and does not vanish
on this one. It spans the one-dimensional first extension
group. Its square is zero because the second group vanishes.
\end{proof}

Allow the root $y$ to move on $X$. A \emph{module with a moving
root} is a left $\mathcal D_X$-module $W$ with a smooth,
$\mathcal O_X$-linear Heisenberg action satisfying
\begin{equation}\label{cbc:rc:moving-relation}
 [\nabla_{\partial_y},J_n]=-nJ_{n-1}.
\end{equation}
For $\mathcal V=\mathcal O_X\otimes V_\kappa$ the connection
is $\partial_y+T$. For any left $\mathcal D_X$-module $N$,
the tensor $N\otimes_{\mathcal O_X}\mathcal V$ has its
Leibniz connection and the modes acting on the second factor.

Define
\[
 \mathfrak z^\ell=\mathcal O_X\otimes C^\bullet(V_\kappa),
 \qquad Tx_n=nx_{n+1},\quad T\xi_n=(n+1)\xi_{n+1},
\]
with connection $\partial_y+T$. The differential commutes
with this connection. As a complex this is
$\mathcal Z^\ell=\mathcal O_X\otimes\mathcal Z$:
$J=x_1$, $c=\xi_0$, and $T^nc/n!=\xi_n$.

\begin{theorem}\label{cbc:rc:moving-endomorphisms}
For every complex $N$ of left $\mathcal D_X$-modules, the
singleton action of Theorem~\ref{cbc:rc:cochain-monoid}
induces a natural equivalence
\begin{equation}\label{cbc:rc:moving-adjunction}
 R\operatorname{Hom}_{\mathcal D_X}(N,\mathcal Z^\ell)
 \simeq
 R\operatorname{Hom}_{\mathrm{mov}}
 (N\otimes_{\mathcal O_X}^{\mathbf L}\mathcal V,\mathcal V).
\end{equation}
The right-hand mapping complex is formed among modules with
a moving root. The tensor with $\mathcal V$ is already exact.
The derivation $K\xi_n=\kappa^{-1}x_{n+1}$, $Kx_n=0$,
satisfies $dK+Kd=T$. Hence the cohomology of $\mathcal Z^\ell$
is $\mathcal O_X\otimes\Lambda(\xi_0)$ with trivial connection.
\end{theorem}
\begin{proof}
The invariant module $W^{\mathfrak b}$ is preserved by the
connection, by~\eqref{cbc:rc:moving-relation}. Induction from
the vacuum gives an adjunction
\[
 \operatorname{Hom}_{\mathrm{mov}}
 (N\otimes_{\mathcal O_X}\mathcal V,W)
 =\operatorname{Hom}_{\mathcal D_X}(N,W^{\mathfrak b}).
\]
Its left adjoint is exact, so the right adjoint preserves
injectives. The continuous cochains compute its derived
functor.

To verify this last assertion with the connection
present, put $A_\kappa=U(\mathfrak h_\kappa)/(\mathbf c-1)$.
The algebra obtained by adjoining $y,D$ with
\[
 [D,y]=1,\quad [y,J_n]=0,\quad [D,J_n]=-nJ_{n-1}
\]
is free as a right $A_\kappa$-module on $y^aD^r$.
The derivation of the modes preserves their relations, since
its value on the relation for $J_m,J_n$ is
$-\kappa m(m+n-1)\delta_{m+n-1,0}=0$.
Induction to this algebra is exact and preserves smoothness:
\[
 J_nD^r=\sum_{j=0}^r\binom rj
 n(n-1)\cdots(n-j+1)D^{r-j}J_{n-j}.
\]
If all $J_n$ with $n\geq p$ kill a vector, all $J_n$ with
$n\geq p+r$ kill its induced $D^r$-multiple.

Thus forgetting the moving connection preserves injectives.
Lemma~\ref{cbc:rc:smooth-ext} proves cochain acyclicity on
them, also as differential modules. The first-quadrant
double complex of an injective resolution now identifies
$R\Gamma_{\mathfrak b}\mathcal V$ with $\mathfrak z^\ell$.
Derived adjunction gives~\eqref{cbc:rc:moving-adjunction}
for all $N$, with derived mapping complexes understood via
injective resolutions and their product totalizations.

It remains to identify this adjunction with the actual
singleton action. Put $\mathcal W=\mathcal O_X\otimes W_1$
and let the current act by the modes of $K_1$.
The action is the roof of $i_1$ and
$z\otimes v\mapsto j_1(z)_{(-1)}i_1(v)$.
Under the underived adjunction this map is $z\mapsto j_1(z)$.

The cochains of $W_1$ have differential
$Q+(:cK_1:)_{(0)}$. Its two maps from $\mathcal Z$ are
\[
 f_0:J\mapsto K_1, c\mapsto c,
 \qquad f_1:J\mapsto J-\kappa Tt, c\mapsto-\eta.
\]
The map fixing $J,c$ is joined to them by
\begin{align*}
 J&\longmapsto J+sp+ds\,\iota,&c&\longmapsto c,\\
 J&\longmapsto J-s\kappa Tt,&c&\longmapsto(1-s)c-s\eta-ds\,t.
\end{align*}
Both are dg vertex maps.

Here the cochain differential sends
$t$ to $\eta+c$, sends $J$ to $\kappa Tc$, and kills
$p,\eta,c$. Its value on $\iota$ is $p$.
Integrating the coefficient of $ds$ gives a chain homotopy
commuting with $T$.

The inclusion of the commutant into these cochains is a
quasi-isomorphism. In the coordinates of
Proposition~\ref{cbc:rc:tower-strict}, the modes of $K_1$
differentiate the variables $\tau_n^{(1)}$. Their continuous
cochains contract the polynomial de Rham pairs. This
contraction is valid for the total complex: give $J,p,\iota$
weight one, $t,\eta$ weight zero, and let $T$ raise weight.

At each weight only finitely many positive-index ghosts can
occur, with at most one zero-index ghost. The cochain-degree
filtration is therefore finite there. Taking the union of
weights preserves the equivalence. The maps $f_0,f_1$ now
show that the singleton roof corresponds to the identity of
$\mathfrak z^\ell$ under derived adjunction.

Finally $dK+Kd=T$ follows on $x_n$ and $\xi_n$ from
$dx_n=\kappa n\xi_n$ and $d\xi_n=0$. Both sides are
derivations, so the equality holds on all states.
\end{proof}

The zero-mode class has disappeared if the allowed modules
are required to have $J_0=0$. At a fixed nonzero charge
$\mu$, the cochain differential instead contains
$\mu\xi_0\wedge$, contracted by $\mu^{-1}\partial_{\xi_0}$.
Thus the extension class belongs to the stated category of
smooth modules with unrestricted zero mode. Its singleton
universal property does not assert a universal property on
the full collision atlas.

\section{The collision class of the product}
\label{cbc:rc:collision-class}

The complex of cochains has only two cohomology classes, but
its product distinguishes the input orders. To detect
that order after applying the coalgebra functor, retain the
homotopy between the binary map and its restriction to a cut.
It has a scalar component in the relative de Rham cohomology
of the normal punctured line.

\subsection{A constant coefficient model}

Let $\mathcal K$ be the regular commutative dg vertex algebra
freely generated by $h$ of degree zero and $c$ of degree one,
with $dh=Tc$, $dc=0$. All its brackets vanish. Put
$E=\Lambda(\varepsilon)$, where $|\varepsilon|=1$,
$T\varepsilon=0$, and the differential is zero.
The map $q:\mathcal K\to E$ sends $c$ to $\varepsilon$
and $h$ to zero.

\begin{proposition}\label{cbc:rc:constant-cochains}
There is a zigzag of quasi-isomorphisms of dg vertex algebras
\[
 \mathcal Z\xrightarrow{j_1}\mathcal Z^{(1)}
 \xleftarrow{g}\mathcal K\xrightarrow{q}E,
 \qquad
 g(h)=\kappa^{-1}(J-\kappa Tt+\tfrac12p),\quad g(c)=-\eta.
\]
It induces equivalences of pointed factorizing coalgebras.
No strict dg vertex map $\mathcal K\to\mathcal Z$ induces
an isomorphism on cohomology.
\end{proposition}
\begin{proof}
Write $h_n=T^nh/n!$ and $c_n=T^nc/n!$.
Then $dh_n=(n+1)c_{n+1}$. The Euler contraction on the
pairs $(h_n,c_{n+1})$ leaves exactly $\Lambda(c_0)$,
so $q$ is a quasi-isomorphism.

Put $U=J-\kappa Tt$. The brackets
\[
 [U_\lambda U]=\kappa\lambda,
 \qquad [p_\lambda U]=[U_\lambda p]=-\kappa\lambda
\]
give $[(U+\tfrac12p)_\lambda(U+\tfrac12p)]=0$.
The image of $c$ commutes with this field, and
$Qg(h)=-T\eta=Tg(c)$. Both images commute with $K_1$.
Thus $g$ is a dg vertex map. It carries the surviving
class $[c_0]$ to $[-\eta_1]$, so
Proposition~\ref{cbc:rc:tower-strict} makes it a
quasi-isomorphism.

The residue construction preserves these equivalences.
For $j_1$ and $g$ this is
Lemma~\ref{cbc:va:free-field-tensor}. For the regular
quotient $q$, use the finite partition filtration of
the Cousin complex.

Each derived tensor coefficient
preserves its quasi-isomorphism, and the residue maps
are natural. Filtered comparison proves the result at
every finite configuration, including collisions and
all cuts. The same argument applies to the algebraic
polynomial path resolutions used below. All maps
preserve the vacuum.

Finally suppose $f:\mathcal K\to\mathcal Z$ is strict.
Lemma~\ref{cbc:rc:smooth-ext} gives
$f(c)=ac+dP$ with $P\in V_\kappa$ and $a\in k$.
Since the closed degree-zero states of $\mathcal Z$
are constants,
\[
 f(h)=a\kappa^{-1}J+TP+n\mathbf1\qquad(n\in k).
\]

The coefficient of $\lambda\mathbf1$ in
$[f(h)_\lambda f(h)]$ is $a^2\kappa^{-1}$.
The cross terms have zero such coefficient because
$J_{(0)}P=0$ and translation has no vacuum component.
The term $[TP_\lambda TP]=-\lambda(\lambda+T)[P_\lambda P]$
also has zero coefficient there. But $[h_\lambda h]=0$,
so $a=0$. Thus $f$ kills the surviving ghost class.
\end{proof}

The exact correction $p/2=Q(\iota/2)$ is essential to this
comparison: it cancels the bracket of $U$ with itself.
The zigzag identifies the pointed object with $C_E$.
Transport the tower product along this zigzag and denote
it by $P_{\mathrm t}:C_E\otimes^!C_E\to C_E$.
Also let $P_{\mathrm c}$ be the product induced by the
ordinary exterior multiplication. These will have
different relative classes.

\subsection{The normal relative complex}

Set $B=E\otimes E=\Lambda(a,b)$, with
$a=\varepsilon\otimes1$ and $b=1\otimes\varepsilon$.
The tensor comparison identifies $C_B$ with
$C_E\otimes^!C_E$. Write $\epsilon_E:E\to k$ and
$\epsilon_B:B\to k$ for their augmentations, and
$\sigma:B\to B$ for exchange, so $\sigma(ab)=-ab$.

Choose the ordered chart $Y=X^2$ with coordinates $(z,w)$,
$u=z-w$, and $j:U=\{u\ne0\}\hookrightarrow Y$.
Put $J_Y=\omega_Y[2]$ with density $dz\wedge dw$ and
$J_U=j^!J_Y$. Consider
\begin{equation}\label{cbc:rc:relative-complex}
 \mathcal R=\operatorname{Fib}\left(
 R\operatorname{Hom}_{\mathcal D_Y}(J_Y,J_Y)
 \longrightarrow
 R\operatorname{Hom}_{\mathcal D_U}(J_U,J_U)\right).
\end{equation}
All differential modules here are right modules. In global
algebraic de Rham complexes its convention is
\[
 \mathcal R^n=\Omega_Y^n\oplus\Omega_U^{n-1},
 \qquad D(\alpha,\beta)=(d\alpha,j^*\alpha-d\beta).
\]

\begin{lemma}\label{cbc:rc:relative-line}
The only nonzero cohomology of $\mathcal R$ is
\[
 H^2(\mathcal R)=k\gamma,
 \qquad \gamma=[(0,du/u)].
\]
Thus a map $J_Y[-2]\to J_Y$ with a specified nullhomotopy
on $U$ has a scalar relative class. Its underlying map on
$Y$ has zero class.
\end{lemma}
\begin{proof}
Polynomial integration contracts the de Rham complex of
$Y$ to $k$. On $U=\mathbb A^1_w\times\mathbb G_{m,u}$,
integration in $w$ leaves
$k[u,u^{-1}]\xrightarrow{\partial_u}k[u,u^{-1}]du$.
Every monomial one-form except $u^{-1}du$ is a derivative.
The two spaces of constants agree under restriction.
The long exact sequence of the fiber gives the result.
The absolute class lies in $H^2_{\mathrm{dR}}(Y)=0$.
\end{proof}

In $B\otimes B$, with factors at positions $z,w$, define
\begin{equation}\label{cbc:rc:relative-inputs}
 v_1=a\otimes b-1\otimes ab,
 \qquad v_2=b\otimes a+1\otimes ab.
\end{equation}
They have degree two and zero product. The constant vertex
field is multiplication, so every Laurent multiple of
$v_i$ has zero collision residue. Hence they define maps
$s_i:J_Y[-2]\to C_{B,2}|_Y$ in the Cousin complex.

For a map $P:C_B\to C_E$ under the vacuum, put
$q_P=C_{\epsilon_E}P$. Its cut compatibility supplies
a homotopy on $U$ from $(q_P)_2$ to the external square
of $(q_P)_1$. Choose a pointed singleton homotopy from
$(q_P)_1$ to $\epsilon_B$. Its external square, followed
by the cut homotopy, gives a nullhomotopy of $(q_P)_2s_i$ on $U$.
Let $\lambda_i(P)\gamma$ be the resulting relative class.

\begin{proposition}\label{cbc:rc:relative-invariant}
The construction gives well-defined functions
\[
 (\lambda_1,\lambda_2):
 \pi_0\operatorname{Map}_{J/}(C_B,C_E)\longrightarrow k^2.
\]
They are independent of the chosen singleton homotopy,
and satisfy
\[
 \lambda(P_{\mathrm c})=(0,0),\qquad
 \lambda(PC_\sigma)=(\lambda_2(P),\lambda_1(P)).
\]
The domain retains both ordered cuts.
\end{proposition}
\begin{proof}
Spencer resolution and polynomial integration identify
\[
 R\operatorname{Hom}_{\mathcal D_X}
 (\mathcal O_X\otimes B,\mathcal O_X)
 \simeq\operatorname{Hom}_k(B,k).
\]
Fixing the value on $1$ leaves a connected mapping space,
so the required singleton homotopy exists. Its loops have
classes $\operatorname{Hom}_k(B^1,k)=k^2$.
A representative $h$ has degree minus one and can be
nonzero only on $a,b$.

The first variation of its external
square at $\epsilon_B$ is
$h\boxtimes\epsilon_B+\epsilon_B\boxtimes h$, with tensor
Koszul signs. Both terms vanish on $v_i$: the other odd
input is killed by $\epsilon_B$, and $h(ab)=0$.
Equivalently, represent the loop by sending $a,b$ to
multiples of $ds$. Its square vanishes because $ds^2=0$.
The relative class therefore does not depend on this
choice.

Evaluation on a chart and its cut preserves homotopies
of full coalgebra maps. The relative class consequently
depends only on the component of $P$.
For $P_{\mathrm c}$, $q_P$ is the strict augmentation,
which kills $v_i$ with constant cut homotopy. Exchange
sends $v_1$ to $v_2$ and $v_2$ to $v_1$, proving the
last formula.
\end{proof}

\subsection{Computing the class on a resolved map}

We compute the invariant through an algebraic path
resolution. Suppose a derived character $f:B\to k$ is
presented by a roof
$B\xleftarrow{\pi}S\xrightarrow{\varphi}k$, where $\pi$
is a surjective quasi-isomorphism of dg vertex algebras.
Put $I=\ker\pi$, an acyclic ideal. Choose closed lifts
$\tilde a,\tilde b\in S^1$. Since $Ta=Tb=0$, choose
$y_a,y_b\in I^0$ with
$Qy_a=T\tilde a$, $Qy_b=T\tilde b$.

The zero mode of $y_a$ commutes with $Q$, because
$(T\tilde a)_{(0)}=0$. Thus one can choose $x_1,x_2\in I^0$
with
\begin{equation}\label{cbc:rc:character-bracket}
 Qx_1=(y_a)_{(0)}\tilde b,
 \qquad Qx_2=(y_b)_{(0)}\tilde a.
\end{equation}

The values $\varphi(x_i)$ do not depend on these choices.
Changing a primitive by $Qe$ adds a bracket whose image under
$\varphi$ is zero. Changing a closed lift by $Qw$ changes its
translation primitive by $Tw$, whose zero mode vanishes. Two remaining choices of $x_i$
differ by a boundary in $I$.

\begin{lemma}\label{cbc:rc:character-reduction}
If $C(f)=C_{\epsilon_E}P$, then
\[
 \lambda(P)=(\varphi(x_1),\varphi(x_2)).
\]
\end{lemma}
\begin{proof}
Write states to the left of differential forms. On two
positions the Spencer differential is
\begin{align*}
 \mathsf D(x\otimes y\,\zeta)
  ={}&Q(x\otimes y)\,\zeta\\
  &+(-1)^{|x|+|y|}\bigl(Tx\otimes y\,dz\wedge\zeta
       +x\otimes Ty\,dw\wedge\zeta+x\otimes y\,d\zeta\bigr).
\end{align*}
The two-position Cousin complex has open and diagonal
summands with differential
$(\sigma,\delta)\mapsto(Q\sigma,r_{12}\sigma-Q\delta)$.
A closed Spencer element is therefore a pair
$(\Sigma,\Gamma)$ with
$\mathsf D\Sigma=0$ and $r_{12}\Sigma=\mathsf D\Gamma$.
On constant coefficients its residue is
\[
 r_{12}(x\otimes y)=
       \sum_{n\geq0}x_{(n)}y\,[u^{-n-1}],\qquad dz=du+dw.
\]

Choose $y_{ab}\in I^1$ with
$Qy_{ab}=-T(:\tilde a\tilde b:)$, and put
$\Xi_{ab}=: \tilde a\tilde b :+y_{ab}\,dw$.
The element
\begin{align*}
 \Sigma_1={}&\tilde a\otimes\tilde b
 +\tilde a\otimes y_b\,dw-y_a\otimes\tilde b\,dz\\
 &+y_a\otimes y_b\,dz\wedge dw-1\otimes\Xi_{ab}
\end{align*}
is closed and lifts $v_1$. Its residue has coefficients
in $I$, whose Spencer complex is acyclic. Choose
$\Gamma$ there with $\mathsf D\Gamma=r_{12}\Sigma_1$.
Then $(\Sigma_1,\Gamma)$ lifts $s_1$.

A singleton homotopy from $\varphi$ to $\epsilon_B\pi$
gives, by its external square, a regular form $\beta$
on $Y$ with $\mathsf D\beta=\varphi\Sigma_1$.
Its restriction to $U$ is precisely the nullhomotopy
in Proposition~\ref{cbc:rc:relative-invariant}.
Subtract the relative boundary of $((\beta,0),0)$.
Since a regular scalar form has zero collision residue,
the class becomes $((0,\varphi\Gamma),0)$.
Under this Cousin resolution, $\gamma=(0,du/u)$ is
represented by $((0,-[u^{-1}]du),0)$.

The principal-part de Rham complex is
\[
 [u^{-m}]\longmapsto-m[u^{-m-1}]du\qquad(m\geq1),
\]
tensored with the polynomial de Rham complex in $w$.
Its first cohomology is one-dimensional, detected by
the coefficient of $[u^{-1}]du$.

If $g\in I^0$ is
that coefficient in $\Gamma$, the same coefficient
in $\mathsf D\Gamma$ is $Qg$: differentiating a
principal part cannot produce a simple pole.
In $r_{12}\Sigma_1$ it is $-(y_a)_{(0)}\tilde b$.
Consequently $g+x_1$ is closed in $I$ and hence exact.
The scalar residue of $\varphi\Gamma$ is
$\varphi(g)=-\varphi(x_1)$.

The minus sign in the representative of $\gamma$
gives $\lambda_1(P)=\varphi(x_1)$.
Exchange $a,b$ for the second assertion. The added
vacuum term in $v_2$ has no residue.
\end{proof}

\subsection{The value for every input order}

For $1\leq\alpha\leq n$ put
\[
 h_\alpha=\kappa^{-1}(U_\alpha+\tfrac12p_\alpha),
 \qquad c_\alpha=-\eta_\alpha,
 \qquad
 g_{n,\alpha}:\mathcal K\to\mathcal Z^{(n)},\quad
 (h,c)\longmapsto(h_\alpha,c_\alpha).
\]
Each is a quasi-isomorphism by the same half-correction
calculation as Proposition~\ref{cbc:rc:constant-cochains}.
For different indices their images commute. Hence their
joint map is a dg vertex map
$A_n:\mathcal K^{\otimes n}\to\mathcal Z^{(n)}$.

The maps $g_{n,\alpha}$ are polynomially homotopic. For
$\alpha>1$ put
\[
 \theta_\alpha=t_1-t_\alpha,\qquad
 w_\alpha=\kappa^{-1}
 (\tfrac12\iota_1+\sum_{1<l<\alpha}\iota_l+\tfrac12\iota_\alpha).
\]
With parameter forms on the left, the homotopy is
\begin{align*}
 h&\longmapsto(1-s)h_1+sh_\alpha+ds\,w_\alpha,\\
 c&\longmapsto(1-s)c_1+sc_\alpha+ds\,\theta_\alpha.
\end{align*}
Indeed $Q\theta_\alpha=c_\alpha-c_1$ and
$Qw_\alpha=h_\alpha-h_1-T\theta_\alpha$.
The only potentially nonzero bracket between these
images has $ds$ coefficient
$(1-s)/(2\kappa)+s/(2\kappa)-1/(2\kappa)=0$.

Transport the $n$-fold tower product to $C_E$ and call
it $P_{\mathrm t,n}$. Its composite with augmentation
is induced by the derived character
\begin{equation}\label{cbc:rc:tower-character}
 \epsilon_{\mathcal K}\circ g_{n,1}^{-1}A_n(q^{-1})^{\otimes n}.
\end{equation}
To see this, include the first quartet by
$\operatorname{in}_n:\mathcal Z^{(1)}\to\mathcal Z^{(n)}$.
This inclusion and the merge agree on $j_1(\mathcal Z)$,
and $\operatorname{in}_ng=g_{n,1}$.
The iterated product carries $g^{\otimes n}$ to $A_n$.
Substituting these identities in the transported roof
proves~\eqref{cbc:rc:tower-character}.

\begin{theorem}\label{cbc:rc:ordered-product-class}
The tower product, its opposite, and the ordinary exterior
product have respective relative classes
\[
 \lambda(P_{\mathrm t})=
  \left(-\frac1{2\kappa},\frac1{2\kappa}\right),\qquad
 \lambda(P_{\mathrm t}C_\sigma)=
  \left(\frac1{2\kappa},-\frac1{2\kappa}\right),\qquad
 \lambda(P_{\mathrm c})=(0,0).
\]
They belong to three distinct components of
$\operatorname{Map}_{J/}(C_E\otimes^!C_E,C_E)$.
More generally, the $n!$ permutations of the inputs
of $P_{\mathrm t,n}$ are pairwise nonhomotopic.
\end{theorem}
\begin{proof}
Choose distinct $\alpha,\beta$ and restrict
\eqref{cbc:rc:tower-character} to these two inputs.
The preceding homotopy allows $g_{n,1}^{-1}$ to be
replaced by $g_{n,\alpha}^{-1}$.
We construct the roof needed in
Lemma~\ref{cbc:rc:character-reduction}.

Let $P$ consist of pairs $(k_0,\gamma)$, where
$k_0\in\mathcal K$ and
$\gamma\in\mathcal Z^{(n)}\otimes k[s,ds]$ satisfy
$\gamma(0)=g_{n,\alpha}(k_0)$.
Projection $P\to\mathcal K$ is a quasi-isomorphism,
since the kernel of evaluation at $0$ on polynomial
forms is contracted by integration. Evaluation at $1$
is surjective: $(0,sz)$ lifts every homogeneous $z$.
It is a quasi-isomorphism because its homotopy to
evaluation at $0$ compares it with $g_{n,\alpha}$.

Take the pullback $S$ with $\mathcal K\otimes\mathcal K$
along $g_{n,\alpha}\otimes g_{n,\beta}$.
Then its projection, followed by $q\otimes q$, is a
surjective quasi-isomorphism $\pi:S\to B$.
The character is $\varphi(x,k_0,\gamma)=\epsilon_{\mathcal K}(k_0)$.
The tautological path verifies that this roof presents
the specified derived character.

Closed lifts and a translation primitive are
\[
 \tilde a=(c^{(1)},c,c_\alpha),\qquad
 y_a=(h^{(1)},h,h_\alpha),\qquad
 \tilde b=(c^{(2)},c,\gamma_b),
\]
where
$\gamma_b=(1-s)c_\alpha+sc_\beta+ds(t_\alpha-t_\beta)$.
The element $y_a$ lies in $\ker\pi$ and satisfies
$Qy_a=T\tilde a$.

The first two components of $(y_a)_{(0)}\tilde b$ vanish.
Its third component is $ds\,\mu$, where
\[
 \mu=(h_\alpha)_{(0)}(t_\alpha-t_\beta)
   =\kappa^{-1}\bigl(\tfrac12-[\beta<\alpha]\bigr).
\]
Here $[\beta<\alpha]$ equals one or zero according
to the condition, and $(p_l)_{(0)}t_m=\delta_{lm}$.
The element $x_1=(0,-\mu,(s-1)\mu)$ belongs to
$\ker\pi$ and satisfies
$Qx_1=(y_a)_{(0)}\tilde b$.
Thus $\varphi(x_1)=-\mu$.

Using the roof based on $g_{n,\beta}$ and exchanging
the inputs gives the other coordinate. The pair is
\begin{equation}\label{cbc:rc:ordered-pair-value}
 \kappa^{-1}
 \bigl([\beta<\alpha]-\tfrac12,
       [\alpha<\beta]-\tfrac12\bigr).
\end{equation}
Lemma~\ref{cbc:rc:character-reduction} identifies it
with the relative class.

For $n=2$, the ordered positions are $\alpha=1$,
$\beta=2$. Formula~\eqref{cbc:rc:ordered-pair-value}
and the exchange law prove the three values.
They are distinct because the characteristic is zero
and $\kappa\ne0$.
In any arity, insert these two inputs and vacua in
all other factors. The resulting pair tells which
of the two selected inputs comes first. The relative
orders of all pairs determine the permutation,
proving the last assertion.
\end{proof}

\begin{proposition}\label{cbc:rc:disjoint-products}
After restriction to pairwise distinct configurations,
$P_{\mathrm t}$, $P_{\mathrm t}C_\sigma$ and $P_{\mathrm c}$
are homotopic with all cuts and vacua.
\end{proposition}
\begin{proof}
On the configuration open only the partition into
singletons survives in the Cousin complex. For a
pointed singleton left differential module $M$, this
restriction is
$J_I\otimes M^{\boxtimes I}|_{U_I}$, with the cut
maps given by regrouping factors. This is a functor
of the pointed derived singleton, since external
tensor over $k$ and open restriction preserve
quasi-isomorphisms.

Polynomial de Rham contraction identifies the
degree-zero singleton maps from $B$ to $E$ with
$\operatorname{Hom}^0_k(B,E)$.
All three products send $a,b$ to $\varepsilon$,
$ab$ to zero, and $1$ to $1$. Their singleton maps
are homotopic. Applying the external-power functor
gives the asserted homotopy on every configuration
open, with all its cuts and vacua.
\end{proof}

The half in~\eqref{cbc:rc:ordered-pair-value} comes
from the exact correction that made the constant
model commute. The integer records the input order.
Both are detected by the compatibility at a collision,
although the products agree on the configuration
opens. A further comparison with a universal action
space must retain this compatibility.

\section{A map invisible to natural actions}
\label{cbc:rc:invisible-map}

The next distinction concerns maps between actors. Their
formal coordinate substitution does not record which root
carries an odd output. We construct a map with zero
substitution in every weight and a nonzero antisymmetric
coefficient on two roots.

For a finite-dimensional vector space $W$ put
\[
 E_W=\Lambda(W[-1]),\qquad
 F_W=\operatorname{CE}_\star(\mathfrak l\otimes E_W),
 \qquad A_W=U_\star(\mathfrak l\otimes E_W),\qquad
 \mathfrak l=L_k.
\]
The exterior generators have degree one and zero translation.
Write $F_r$ when a basis identifies $W$ with $k^r$.
These are the constant exterior field coalgebras, with
pointing from and augmentation to $J=F_0$.
All actor maps constructed in this section preserve both.
The suspension $s$ has degree minus one, so a pure input
$s(em)$, with $e\in W[-1]$ and $m\in\mathfrak l$, has
degree zero.

\subsection{Derivations of exterior coefficients}

Let $D$ be an endomorphism of a finite-dimensional space
$V$, extended as an even derivation of $E_V$, of
$\mathfrak g_V=\mathfrak l\otimes E_V$, and of $A_V$.
It fixes $\mathfrak l$. Introduce another exterior generator
$\varepsilon$, placed on the left. For homogeneous $x,m$,
\[
 [x,\varepsilon m]=(-1)^{|x|}\varepsilon[x,m],\qquad
 [\varepsilon m,\varepsilon n]=0.
\]

The tree set $\mathcal T_w$ used in
\eqref{cbc:va:bar-tree} consists of nonnegative sequences
$(a,b_1,\ldots,b_w)$ with
\[
 a+\sum_{i=1}^w b_i=w,\qquad
 a+\sum_{i=1}^j b_i\geq j+1\quad(0\leq j<w).
\]
For homogeneous actor coefficients $m_i$, let
$\chi(\pi;m)$ be their Koszul permutation sign and set
\begin{equation}\label{cbc:rc:graded-tree}
 \theta_w^D(x;m_1,\ldots,m_w)=
 \sum_{(a,b)\in\mathcal T_w}\sum_{\pi\in S_w}
 \frac{\chi(\pi;m)}{a!b_1!\cdots b_w!}
 D^ax\,D^{b_1}m_{\pi(1)}\cdots D^{b_w}m_{\pi(w)}.
\end{equation}
Put $\theta_0^D(x)=x$.
Define a twisting cochain to vanish except on words with
one input from $s\mathfrak g_V$. Move that input first,
with the sign of the suspended permutation, and apply
$\theta_w^D$ to the remaining actor inputs.
The two signs use different degrees:
$|sx|=|x|-1$, whereas $|s(\varepsilon m_i)|=|m_i|$.

\begin{lemma}\label{cbc:rc:graded-action}
This cochain defines a map $\alpha_D:F_{k\oplus V}\to F_V$
whose restriction to $F_V$ is the identity.
On pure inputs, for $z\in V[-1]$ and $n,m_i\in\mathfrak l$,
its value is
\begin{equation}\label{cbc:rc:graded-pure}
 (D^wz)n\,\operatorname{pbw}(m_1\cdots m_w),
 \qquad
 \operatorname{pbw}(m_1\cdots m_w)=
 \frac1{w!}\sum_{\pi\in S_w}m_{\pi(1)}\cdots m_{\pi(w)}.
\end{equation}
It vanishes on pure inputs with any other number of
factors from $V$.
\end{lemma}
\begin{proof}
In the twisting equation a nonzero term has exactly two
inputs from $s\mathfrak g_V$. A source bracket either
joins these or joins one to an actor input. A product
of two twisting values also needs exactly two such
inputs. Every other component vanishes.

On the remaining components use the degree-zero tree
identity proved in Theorem~\ref{cbc:va:profiles-theorem}.
For an input of coefficient degree $d\geq0$, adjoin
$d$ distinct auxiliary exterior variables of degree
minus one and multiply the input by their ordered
product. Use disjoint variables for different inputs.
The new coefficient inputs have degree zero, and $D$
fixes the auxiliary algebra.

Apply the tree identity
there and extract the product of all auxiliary
variables. Permuting the actor inputs gives precisely
$\chi(\pi;m)$, and moving the module input gives
the suspended sign. The extracted monomial is nonzero.
Thus the identity holds with
$Q_2(sx\,sy)=(-1)^{|x|}s[x,y]$ and the convolution
tensor sign $(-1)^{|c'|}$ in every degree.

Lemma~\ref{cbc:va:chevalley-bar} converts the twisting
cochain to a coalgebra map. Every positive-weight
tree has $a\geq1$. Its first factor therefore lies
in the positive exterior ideal, proving compatibility
with augmentation. On the vacuum $D$ vanishes, proving
compatibility with the pointing.

Positivity makes
the bar sum finite at each configuration.
For~\eqref{cbc:rc:graded-pure}, $D$ kills all $m_i$,
so only $(a,b)=(w,0,\ldots,0)$ survives. Its permutation
average is the stated PBW product.
\end{proof}

\subsection{Marked coefficients and composition}

Here natural cochains are universal multilinear expressions in
products, brackets, exterior-coefficient operators and the chosen
derivations, with finite values at each configuration. The bar
comparisons of Section~\ref{cbc:va:profile-recipients} retain this class.
For a natural map $f:F_W\to F_V$, write $f_0:W\to V$
for its linear term.

At two distinct labeled roots,
project the quadratic term to PBW output with one
odd factor. Encode it by a $V$-valued one-form
$\omega_f$ on $W$ with homogeneous linear coefficients.
The term $t_i\,dt_j\,v$ means that the input of color
$j$ carries the odd output $v$ and the input of color
$i$ is unmarked. If $i=j$, sum the two marked-root
possibilities. Thus $t\,dt$ encodes
$(\varepsilon m)n+(\varepsilon n)m$.

These marked coefficients cannot be changed by a homotopy
under and over $J$. Indeed a pure actor input has degree
zero. An output with $b\geq1$ odd factors, together with
$n$ additional ordinary Lie inputs and coefficient
resolution degree $q\geq0$, gives cochain degree at least
$b+n+q$. There is no degree-zero homotopy coefficient
on a pure actor input.

Its remaining variation would
come from a source bracket joining roots. Open
adjunction kills every map from that joined support
to two distinct output roots. The two choices of
marked output are independent PBW summands on the open.

\begin{lemma}\label{cbc:rc:marked-composition}
For composable natural maps,
\[
 (g f)_0=g_0f_0,
 \qquad \omega_{g f}=g_0\omega_f+f_0^*\omega_g,
 \qquad (f_0^*\omega)(t)=\omega(f_0t)\circ f_0.
\]
The first nonlinear form of a diagonal tensor product
is the block sum of its forms.
\end{lemma}
\begin{proof}
Lift the inner twisting cochain to the intermediate
cobar, using the relative bar contraction following
Lemma~\ref{cbc:va:natural-cohomology}.
Write $a_{v,m}=t(s(vm))$ and $M_m=t(sm)$ for its
singleton generators, where $t=s^{-1}$.
On two roots the lift consists of the marked primitive
products prescribed by $\omega_f$ and the binary
generator $t(sx\,sy)$, where
$x=f_0(e_i)m$, $y=f_0(e_j)n$.
On the open its differential is
$-t(sx)t(sy)-t(sy)t(sx)$, exactly the source
coproduct term. Joined-root corrections vanish there.

Applying $g$ to the primitive products gives
$g_0\omega_f$, and applying it to the binary generator
gives $f_0^*\omega_g$.
The marked coefficients are homotopy invariant by
the preceding degree and support argument.

For the tensor assertion, independently grade the
two target exterior factors. The strong tensor
comparison preserves these gradings. Any term
using both factors has at least two odd outputs,
so the one-odd-output component is the block sum.
\end{proof}

The one-variable maps~\eqref{cbc:va:actor-bar-map}
realize any polynomial with zero constant term by
choosing their coefficients. A pointed polynomial
$\Phi:W\to V$ is a finite sum of powers of linear
forms, component by component. The corresponding
parallel one-variable maps, followed by exterior
linear addition, therefore give an actual map
$f_\Phi:F_W\to F_V$. For a homogeneous component $\Phi_r$ of degree $r$, the symmetric
marked form is $r^{-1}d\Phi_r$. Its $r$ input roots contribute the
same coefficient, and their sum is the coordinate polynomial. The independent exterior
gradings exclude cross terms from parallel factors.

Take $V=\langle z_1,z_2\rangle$ with
$Dz_1=z_2$, $Dz_2=0$. Realize
\[
 \Phi(t_1,t_2)=(t_1,t_2,-t_1t_2),
 \qquad
 -t_1t_2=\frac{(t_1-t_2)^2-(t_1+t_2)^2}{4}.
\]
Let $e_1,e_2$ be the source exterior generators dual to $t_1,t_2$.
Relabel its outputs by
$(y_1,y_2,y_3)\mapsto(z_1,\varepsilon,z_2)$,
apply $\alpha_D$, and project $z_1\mapsto0$,
$z_2\mapsto\varepsilon'$ to one target generator.
Call the resulting composite $f:F_2\to F_1$.

\begin{proposition}\label{cbc:rc:antisymmetric-map}
The map $f$ has zero linear term. It admits a twisting representative
with zero pure binary values on equal colors and mixed value
\begin{equation}\label{cbc:rc:antisymmetric-value}
 \tau_f(s(e_1m)s(e_2n))
  =\tfrac12\bigl((\varepsilon'm)n-(\varepsilon'n)m\bigr).
\end{equation}
Its marked form is
$\tfrac12(t_2\,dt_1-t_1\,dt_2)\varepsilon'$,
so it is not homotopic to the map through $J$.
\end{proposition}
\begin{proof}
The last composite from three variables has linear
part $y_3$ and marked form $y_2\,dy_1$.
The direct polynomial realization has linear part
$(t_1,t_2,0)$ and form
$(0,0,-\tfrac12d(t_1t_2))$.
Lemma~\ref{cbc:rc:marked-composition} gives the
stated linear term and antisymmetric form.

For the exact mixed representative, lift its
binary generator in the intermediate cobar by
\[
 L_{12}(m,n)=t(s(z_1m)s(\varepsilon n))
  -\tfrac12\bigl(a_{z_2,m}M_n+a_{z_2,n}M_m\bigr).
\]
The joined source $s(e_1e_2p)$ lifts to
$-t(s(\varepsilon z_1p))$, because
$z_1\varepsilon=-\varepsilon z_1$.
The cobar differential gives
\[
 dL_{12}=-t(s(\varepsilon z_1[m,n]))
         -a_{z_1,m}a_{\varepsilon,n}
         -a_{\varepsilon,n}a_{z_1,m},
\]
which is the required source differential.
The primitive correction is closed.

Applying the last map to the first term gives
$(\varepsilon'm)n$. Its value on the correction
is $-((\varepsilon'm)n+(\varepsilon'n)m)/2$.
Their sum is~\eqref{cbc:rc:antisymmetric-value}.
On either equal-color pair, the inner quadratic
term is zero and the last twisting cochain has
either two module inputs or none, so it vanishes.

The relative bar contraction extends the compatible
lift without changing these binary values.
All constituent maps already exist in every
weight. The marked-form invariant proves the
last assertion.
\end{proof}

\subsection{Vanishing under natural actions}

The natural cohomology in
Lemma~\ref{cbc:va:natural-cohomology} records actor
coordinates and derivation classes. Precomposition
by the maps just constructed preserves this natural
complex. The bar contraction acts on the Lie and
bar factors and retains each exterior coefficient
operator. Thus it lifts the graded tree cochains
and their parameter homotopies as well. The
corrections increase primitive input degree and
are finite at each configuration. They preserve
the pointing, augmentation and actor faces.

On the coordinate algebra the map of
Lemma~\ref{cbc:rc:graded-action} is
\begin{equation}\label{cbc:rc:graded-coordinate}
 (u,z)\longmapsto e^{uD}z.
\end{equation}
To verify the normalization, let $\ell\in V^*$
extract one exterior coefficient on bar length one.
On~\eqref{cbc:rc:graded-pure} it gives
$\ell(D^wz)n\operatorname{pbw}(m_1\cdots m_w)$.
The $w$ identical actor extractors sum $w!$
permutations, while PBW averages them.
Its coordinate is therefore
$u^w\ell(D^wz)/w!$.

Projection to natural cohomology kills all other
nonprimitive exterior labels and fixes the
derivation classes. Composition is ordinary
substitution, and a diagonal tensor product
acts independently on its coordinate blocks.

\begin{proposition}\label{cbc:rc:invisible-action}
The map $f$ of Proposition~\ref{cbc:rc:antisymmetric-map}
has zero coordinate substitution in every weight.
Consider the universal class of Lemma~\ref{cbc:va:natural-cohomology}
for the specified commuting derivations. Every positive-weight natural
action in this class pulls back along $f$ to an action naturally
homotopic to the trivial one. In particular, this holds for the
Heisenberg derivation action $\alpha_D$ with
$D=\partial_{x_1}$ from
Theorem~\ref{cbc:va:profiles-theorem}.
\end{proposition}
\begin{proof}
For the last three-variable map, $D^2=0$ on
$\langle z_1,z_2\rangle$, so
\eqref{cbc:rc:graded-coordinate} and the projection
give $y_3+y_1y_2$.
Its composite with $\Phi$ is identically
$-t_1t_2+t_1t_2=0$ in the completed coordinate
algebra, not merely to second order.

A positive-weight natural action has class
$F\in\prod_{q\geq1}\operatorname{Sym}^q(V^*)
\otimes\mathfrak d$, with the chosen derivation
space $\mathfrak d$. The pullback class is
$F(0)=0$. The natural bar comparison and
weightwise cohomology computation following
Lemma~\ref{cbc:va:natural-cohomology} give the
equivalence of the corresponding complete
Maurer--Cartan spaces. Equivalently, solve the
relative path equation successively in positive
actor weight, using the kernel contraction
there. This gives an all-weight pointed path
to the trivial action. Evaluation on each
coefficient family preserves the path, and
all its sums are finite at a fixed configuration.
\end{proof}

The map has retained information even though
all these natural action classes vanish after
pullback. The next argument shows that its
distinction from the vacuum map survives every
receiving pointed factorizing action.

\section{An obstruction for every pointed recipient}
\label{cbc:rc:pointed-obstruction}

The vacuum forces a degree-zero one-root operation
of a factorizing action to be a module map. This
prevents a nonzero derivation action from being
carried entirely by mixed components after its
singleton actor class has vanished.

Retain $C=C_{V_\kappa}$, $\mathfrak g=L_{V_\kappa}$,
$\mathfrak l=L_k$, and set $A=U_\star\mathfrak g$.
Both diagonal families $\mathfrak g,\mathfrak l$
are in degree zero, whereas $C_1=\mathfrak g_1[1]$
and $J_1=\mathfrak l_1[1]$.
Write $\mu_{\mathfrak g}:\mathfrak g\star\mathfrak g\to\mathfrak g$
for the coefficient bracket.
Let $\mathscr A_C^{\mathrm f}$ be the full category
of pointed actions of Section~\ref{cbc:va:pointed}
whose actor coalgebra is factorizing. Its morphisms
retain the specified path between the actions.
No lower cohomological bound is imposed on an actor.

For an object $(Z,u,\gamma)$ put
\[
 B=\Omega Z,
 \qquad \mathcal T_C=(-)\otimes^!C,
 \qquad \Phi=\Omega\gamma:\mathcal T_C(B)\longrightarrow A.
\]
Lemma~\ref{cbc:va:strong-tensor} makes $\mathcal T_C$
exact and strong symmetric monoidal for $\star$.
The pointing gives $\iota:\mathfrak l\to B$, and
its specified path identifies $\Phi\mathcal T_C(\iota)$
with the generating inclusion $\mathfrak g\to A$.

For a degree-zero class $a:\mathfrak l|_Y\to B|_Y$
on an open $Y\subset X$, define its one-root action
\[
 \Psi(a)=\Phi\mathcal T_C(a)\in
          \operatorname{End}(\mathfrak g|_Y).
\]
Here $A_1=\mathfrak g_1$. Closed adjunction on each small-diagonal branch identifies a map
out of $\mathfrak l$ with its singleton map. This also determines
the map on the full diagonal family.

\subsection{The two ordered products}

On the two-root collision scheme $\mathscr R_2$, let
$j:U\hookrightarrow\mathscr R_2$ be the disjoint open. Its complement
$i$ is the union of the two disjoint diagonal branches.
Write $\mu_{12},\mu_{21}$ for the two ordered
products of $B$ on $U$.

\begin{lemma}\label{cbc:rc:open-products}
Both maps
$\mu_{12},\mu_{21}:(B_1\boxtimes B_1)|_U\to B_2|_U$
are equivalences. If $\sigma$ exchanges the
external factors, the exchange
$R_B=\mu_{12}^{-1}\mu_{21}\sigma$ is natural
under algebra maps between such cobars.
For $A=U_\star\mathfrak g$, $R_A$ is the identity.
\end{lemma}
\begin{proof}
On two roots the reduced cobar has one generator
summand $t\overline Z_2$ and two product summands
$M_{12},M_{21}$, each an external product of
$t\overline Z_1$ with itself.
Its differential from the generator to a product
is the corresponding coproduct, with the cobar
sign. Factorization makes both coproducts
equivalences on $U$. The quotient by $M_{12}$
is the cone of the other equivalence, hence
acyclic. Interchange the two orders for the
second assertion. No connectivity is used.

For an algebra map $F:B\to B'$, compatibility
with both products gives
\[
 (F_1\boxtimes F_1)R_B=R_{B'}(F_1\boxtimes F_1)
\]
after composing with $\mu'_{12}$.
In the envelope $A$, the difference of the
two products is the chiral bracket. It is
supported on the diagonal, so its open
exchange is the identity.
\end{proof}

\begin{lemma}\label{cbc:rc:one-root-module}
For every degree-zero one-root class $a$, the
endomorphism $M=\Psi(a)$ satisfies on both
ordered collision branches
\begin{equation}\label{cbc:rc:module-identity}
 M\mu_{\mathfrak g}
 =\mu_{\mathfrak g}(\mathrm{id}\star M)
 =\mu_{\mathfrak g}(M\star\mathrm{id}).
\end{equation}
If $M$ is also a coefficient derivation, then
$M=0$. On a nonempty open of $X$, every such
one-root endomorphism is a constant scalar.
\end{lemma}
\begin{proof}
Work over $Y$ and put $R=\mathfrak l\boxtimes\mathfrak l$
on its two-root scheme. Each product summand of
$\mathfrak l\star\mathfrak l$ is $j_*j^*R$.
Let $\kappa_a=[\iota,a]$ be the difference of
the two ordered products on that source.
The envelope relation and the vacuum path give
\begin{equation}\label{cbc:rc:commutator-image}
 \Phi\mathcal T_C(\kappa_a)
       =\mu_{\mathfrak g}(\mathrm{id}\star M).
\end{equation}

On $U$, put $w=(R_B-\mathrm{id})(\iota\boxtimes a)$
and let $P(w)=\mu_{12}j_*(w)$, using the global
product from the first convolution summand.
The map $D_a=\kappa_a+P(w)$ restricts to zero on
$U$. Moreover $\Phi\mathcal T_C(P(w))=0$.
Indeed $Z\otimes^!C$ is factorizing, and the
strong monoidal comparison identifies its cobar
exchange with that transported by $\mathcal T_C$.
Naturality in Lemma~\ref{cbc:rc:open-products}
then kills the image of $w$, since $R_A=\mathrm{id}$.
The open projection formula carries this zero
map and a representing nullhomotopy through
$j_*$ and multiplication.

Set $Q=\operatorname{cofib}(j_!j^*R\to j_*j^*R)$.
Open adjunction lifts $D_a$ to $Q\to B_2$.
Only the regular holonomic source $R$ requires
$j_!$ here, so the construction allows unbounded
$B_2$. The localization octahedron gives
\begin{equation}\label{cbc:rc:localization-triangle}
 i_*i^*R\longrightarrow Q\longrightarrow i_*i^!R[1].
\end{equation}

On each branch, the right-module shifts are
$i^!R=\mathfrak l[-1]$ and
$i^*R=\mathfrak l[1]$.
The normal-coordinate sequence
$\omega\to\omega(*0)\to\delta_0$
fixes the first shift, and duality fixes the
second. The composite
$j_*j^*R\to Q\to i_*\mathfrak l$ is precisely
the vacuum principal-part bracket.

Closed projection and $i^!C_2=C_1$ identify
the outer terms after $\mathcal T_C$ with
$i_*\mathfrak g[1]$ and $i_*\mathfrak g$.
The envelope $A_2$ lies in the right-module
heart. Its finite PBW filtration has quotients formed by open and closed
direct images of external degree-zero modules. These direct images
are exact on the affine collision charts.
Also $i^!A_2=\mathfrak g$. Consequently
\[
 \operatorname{Hom}(i_*\mathfrak g[1],A_2)
 =\operatorname{Ext}^{-1}(\mathfrak g,\mathfrak g)=0.
\]

After applying $\mathcal T_C$ and $\Phi$,
the map from $Q$ therefore factors on each
branch through an endomorphism $N$ of
$\mathfrak g$. Equation~\eqref{cbc:rc:commutator-image}
becomes
\[
 \mu_{\mathfrak g}(\mathrm{id}\star M)=N\mu_{\mathfrak g}.
\]

Insert the vacuum in the first slot. Its
bracket is the principal-part map
$\mathrm{can}:j_*j^*(\mathfrak l\boxtimes\mathfrak g)
\to i_*\mathfrak g$, so $(N-M)\mathrm{can}=0$.

Precomposition by $\mathrm{can}$ is injective on these degree-zero
maps. Its kernel is an image of
$\operatorname{Hom}((\mathfrak l\boxtimes\mathfrak g)[1],A_2)$.
This group vanishes because both unshifted objects are in the heart. Hence $N=M$.
Skew symmetry gives the other slot identity.

If $M$ is a derivation, subtract its module
identity from its derivation identity. The
remaining slot is zero. Insert the vacuum
there and use the same injectivity to obtain
$M=0$.

For the scalar assertion, the full Laurent
coefficient source in~\eqref{cbc:rc:module-identity}
gives commutation with every current mode,
including creation modes. To extract them,
use the first coordinate as normal derivative,
so no derivative falls on the second-coordinate
matrix of $M$. The zeroth normal coefficient
of the bracket on $(z-w)^{-n}J\boxtimes v$
is $J_{-n}v$ for $n>0$.

Put $c=M\mathbf1$. Creation-mode commutation
gives $M(f)=fc$ on polynomial states.
Positive-mode commutation gives
$\kappa n\partial_{x_n}c=0$ for every $n>0$.
Thus $c$ is a scalar function of the position.
Horizontalness makes its derivative zero,
so it is constant.
\end{proof}

The derivation $D=\partial_{x_1}$ fails the
module identity already on the vacuum:
\[
 (DJ_{-1}-J_{-1}D)\mathbf1=\mathbf1.
\]
Commutation with annihilation modes alone
would not detect this defect.

\subsection{The singleton class cannot vanish}

We need two facts about derived coefficient
maps. The first controls first actor weight.
Write $P_2=\operatorname{Sym}^2_\star\mathfrak g$,
with its adjoint $\mathfrak g$-module structure.
The Chevalley filtration for
$\operatorname{Ext}^1_{\mathfrak g}(\mathfrak g,P_2)$
has entries
\[
 E_1^{n,q}=\operatorname{Ext}^q_{\mathscr D}
 (\Lambda^n_\star\mathfrak g\star\mathfrak g,P_2),
 \qquad n,q\geq0.
\]
Open adjunction kills column $n=0$ in every
degree: its source has one-root support,
whereas $P_2$ is open direct image from two
distinct roots.

In total degree one, only
$(n,q)=(1,0)$ remains. Its differential is
the module cocycle equation. There are no
boundaries, and no higher differential can
enter or leave it. Thus equality of these
extension classes is equality of their
actual two-output cocycles. This argument
also holds over every open $Y\subset X$.

For a coefficient derivation $D$, the cocycle
is $\gamma_D(x;v)=m(Dx,v)$, where $m$ is
symmetric multiplication into $P_2$.
Its module defect is
$m([x,Dy]+[Dx,y]-D[x,y],v)=0$.
Its open operation is $D\otimes\mathrm{id}$,
so its class is nonzero if $D\ne0$.
This is the first-weight class of the action
$\alpha_D$: its first tree operation is
exactly $Dx\,v$.

The second fact is finite tensor expansion.
For an affine open $Y\subset X$, an arbitrary
complex $N$ of differential modules, and an
arbitrary complex $W$ over $k$,
\begin{equation}\label{cbc:rc:coefficient-kunneth}
 R\operatorname{Hom}_Y(\mathfrak l,N\otimes_k W)
 \simeq R\operatorname{Hom}_Y(\mathfrak l,N)\otimes_k W.
\end{equation}
Under unshifted left/right conversion,
$\mathfrak l$ becomes $\mathcal O_Y$.

Its Spencer resolution on the affine coordinate is
\[
 0\longrightarrow D_Y
 \xrightarrow{P\mapsto P\partial_y}D_Y
 \longrightarrow\mathcal O_Y\longrightarrow0.
\]
It is finite free, so its Hom complex commutes
with algebraic tensor by any $W$, including
an unbounded complex. Over $k$, K\"unneth is
an algebraic direct sum. Every class consequently
has a finite expression in its bidegrees.

For a point $q\in X(k)$ normalize a right
differential-module fiber by $N_q=i_q^!N[1]$.
Then $\mathfrak l_q=k$ and $\mathfrak g_q=V_\kappa$
are in degree zero.
Let $\Psi_q$ denote the action on these normalized fibers.
Given a morphism
$h:(F_1,\alpha_D)\to(Z,\gamma)$ in
$\mathscr A_C^{\mathrm f}$, let $h_0$ be its
twisting component on one pure actor input.
Its plain singleton class is
\[
 [h_0]\in H^1R\operatorname{Hom}(\mathfrak l,B).
\]
This class forgets equivariance by $\mathfrak l$.

\begin{theorem}\label{cbc:rc:singleton-survival}
For the action of a nonzero coefficient derivation
$D$, every receiving morphism $h$ has
$[h_0]|_Y\ne0$ on every nonempty affine open
$Y\subset X$.
\end{theorem}
\begin{proof}
Suppose this plain class vanishes. Resolve the
pointed arrow before forming its relative
convolution complex. Its bar filtration gives
a cofibrant graded quotient, so restriction
of the resolved Hom complexes is surjective.
Fixing the twisting cochain on the resolved
vacuum therefore computes the homotopy fiber
of restriction with its specified path.

A primitive of $h_0$ extends to a degree-zero
relative cochain, zero on the vacuum. Gauge
by its exponential in positive actor weight.
At one pure input this changes $h_0$ by its
differential and sets it to zero exactly.
Higher terms affect only higher input number.
All inverse and exponential series are finite
at each configuration, since every positive
factor consumes a nonempty set of labels.
This gives $h'\simeq h$ with $h'_0=0$, while
retaining the action-compatibility path.

The first-weight term of $h'$ begins with a
closed degree-zero component $c_1$ on one
ordinary input and one actor input. There
is no support restriction on $c_1$.
In the normalized bar map of
$\mathcal T_C(h')$, its binary component
is $\mathcal T_C(c_1)$ plus products involving
$\mathcal T_C(h'_0)$. The products vanish.
After the action the component is therefore
$\Phi\mathcal T_C(c_1)$.

Choose $q\in Y(k)$ and put $Y'=Y\setminus\{q\}$.
Such a point exists because $k$ is infinite.
Place the ordinary input at the varying root
in $Y'$ and the actor input at $q$.
Factorization in the first product order
identifies the target with $B|_{Y'}\otimes B_q$.

Formula~\eqref{cbc:rc:coefficient-kunneth}
expresses $[c_1]$ as a finite sum of bidegrees
$(r,-r)$. After applying the action, the
second factor maps to
$H^{-r}R\operatorname{Hom}_k(V_\kappa,V_\kappa)$,
which vanishes for $r\ne0$.
Thus its open two-output operation has the form
\[
 \sum_{i=1}^m M_i\otimes N_i,
 \qquad M_i=\Psi(a_i),\quad N_i=\Psi_q(b_i),
\]
with $a_i,b_i$ degree-zero classes.

The specified action path identifies its
first-weight extension class with that of
$\alpha_D$. The preceding Chevalley filtration
has no boundaries in its $P_2$ component.
Hence on this slice the equality is an
equality of coefficient operations:
\[
 D\otimes\mathrm{id}_{V_\kappa}
        =\sum_{i=1}^m M_i\otimes N_i.
\]

Algebraic tensors of these endomorphisms inject
into external endomorphisms. For a finite
independent list in the second factor,
finitely many vector evaluations and linear
functionals separate its members.
Choose a linear functional $\lambda$ on the
span of $\mathrm{id},N_1,\ldots,N_m$ with
$\lambda(\mathrm{id})=1$. It gives
$D|_{Y'}=\sum_i\lambda(N_i)M_i$.

By Lemma~\ref{cbc:rc:one-root-module}, every
$M_i$ is a module map. Thus $D|_{Y'}$ is both
a derivation and a module map, and is zero.
The state differential module is torsion-free,
and the image of each state is algebraic.
Localization is faithful on that image, so
$D=0$ on $X$, a contradiction.
\end{proof}

\subsection{The parallel pair survives}

Let $f:F_2\to F_1$ be the map of
Proposition~\ref{cbc:rc:antisymmetric-map}, and
let $f_0$ factor through $J$.
For $D=\partial_{x_1}$,
Proposition~\ref{cbc:rc:invisible-action} supplies
an all-weight pointed path $p_f$ from the
pullback of $\alpha_D$ along $f$ to the
trivial action. Use the constant path $p_0$
for $f_0$. They give two points
\begin{equation}\label{cbc:rc:pointed-pair}
 (f,p_f),(f_0,p_0)\in
 \operatorname{Map}_{\mathscr A_C^{\mathrm f}}
 ((F_2,\mathrm{triv}),(F_1,\alpha_D)).
\end{equation}

We first specify the invariant used to separate
their images in a general recipient. Suppose two maps $F_2\to Z$ under $J$ have zero pure
single-actor components. A pointed homotopy between them preserves
the class of their pure binary component in
\begin{equation}\label{cbc:rc:binary-obstruction}
 H^1R\operatorname{Hom}
 (j^*(\mathfrak l\star\mathfrak l),j^*B),
\end{equation}
with both input labels retained.

To prove this,
use the full resolved coefficient differential
$\partial$ and write the interval homotopy
as $\tau(u)+c(u)du$.
Its two singleton components have the form
$\partial E(u),\partial F(u)$, with
$E(0)=F(0)=0$ and
$\partial E(1)=\partial F(1)=0$.
They are obtained by integrating the singleton
components of $c$. On the disjoint open the
source bracket vanishes, and the binary
interval equation reads
\[
 \tau'_{1,1}=\partial c_{1,1}
       +[\partial E,F']+[\partial F,E'].
\]

The full Leibniz rule, with $|E|=|F|=0$, gives
\[
 [\partial E,F']+[\partial F,E']
   =\partial[E,F']-\frac{d}{du}[E,\partial F].
\]
The integrated endpoint term vanishes.
Thus the change of binary components is a
boundary in~\eqref{cbc:rc:binary-obstruction}.
The commutator cochains need not vanish.

\begin{theorem}\label{cbc:rc:pointed-nonrepresentability}
For every morphism
$h:(F_1,\alpha_D)\to(Z,\gamma)$ in
$\mathscr A_C^{\mathrm f}$, the maps $hf$ and
$hf_0$ are not homotopic under $J$.
Consequently the category of pointed factorizing
actions is not filtered, and
\[
 \Phi_C(F,u)=\operatorname{Map}_{C/}(F\otimes^!C,C)
\]
is not representable on pointed factorizing
coalgebras.
\end{theorem}
\begin{proof}
Both composites have zero pure singleton
components. In the source cobar, on distinct
roots, the signs give
\[
 d\,t(s(\varepsilon n)s(m))
   =-t(s(\varepsilon n))t(s(m))
      +t(s(m))t(s(\varepsilon n)).
\]
After $\Omega h$, this replaces the reversed
product $h_0(n)\iota(m)$ by $\iota(m)h_0(n)$
modulo a boundary. This identity comes from
the cocommutative source and imposes no
centrality on the recipient's vacuum.

Fix $q\in X(k)$ and put $Y=X\setminus\{q\}$.
Let the root of $m$ vary in $Y$ and the
root of $n$ be $q$. In the first product
order, the binary class of $hf$ is
\[
 \tfrac12\bigl([h_0]|_Y\otimes[\iota_q]
                    -[\iota]|_Y\otimes[h_0]_q\bigr).
\]

By~\eqref{cbc:rc:coefficient-kunneth}, its two
terms lie in distinct bidegrees $(1,0)$ and
$(0,1)$. The first is nonzero:
Theorem~\ref{cbc:rc:singleton-survival} gives
$[h_0]|_Y\ne0$, and
$\Psi_q(\iota_q)=\mathrm{id}_{V_\kappa}$ gives
$[\iota_q]\ne0$.
No claim about the second term is needed.
The binary class of $hf_0$ is zero.

Normalized restriction to $Y\times\{q\}$
is a map of derived coefficient complexes.
It preserves binary components and boundaries,
so these distinct slice classes come from
distinct classes in~\eqref{cbc:rc:binary-obstruction}.
The homotopy invariant just proved excludes
a homotopy under $J$, and hence excludes one
in the action category with its additional
action two-cell.

A filtered infinity category equalizes every
parallel pair after some receiving map.
The pair~\eqref{cbc:rc:pointed-pair} violates
this requirement. If $\Phi_C$ were represented,
its universal action would be a final object
of $\mathscr A_C^{\mathrm f}$. Its mapping
space from $(F_2,\mathrm{triv})$ would be
contractible and would identify the two
composites, again a contradiction.
\end{proof}

The two actor coalgebras admit additive exterior
monoid structures. The map $f$ is not a map
of those monoids: its first form
$\tfrac12(t_2\,dt_1-t_1\,dt_2)$ has a nonzero
cross term under addition. Therefore
Theorem~\ref{cbc:rc:pointed-nonrepresentability}
does not rule out a final transverse monoid
action. The specific cochain monoid is
already nonfinal by
Theorem~\ref{cbc:va:cochain-nonfinal}.
Constructing another final monoid, or proving
an obstruction within the category of monoid
actions and monoid morphisms, remains a
different universal-property problem.

\section{Field reconstruction and vacuum normalization}
\label{cbc:rc:field-reconstruction}

A reconstruction from collision coefficients must
recover the fields, their exchange, and their vacuum
normalization together. Binary field evaluation can
lose an exchange operator. An outer collision with
a vacuum partner recovers it while retaining the
intervening state.

Let $\mathsf R$ be a commutative algebra over a
characteristic-zero field. Let $V$ be a flat
$\mathsf R$-module, concentrated in degree zero,
with $\mathsf R$-linear translation $T$ and a vacuum $\mathbf1$
spanning a free direct summand. Suppose $\mathsf R$-bilinear fields
$Y(a,u)b\in V((u))$ satisfy creation and translation,
\[
 Y(\mathbf1,u)=1,\quad Y(a,u)\mathbf1=e^{uT}a,
 \quad [T,Y(a,u)]=\partial_uY(a,u)=Y(Ta,u).
\]
Also assume weak associativity: for each triple, after
multiplication by a sufficiently high power of
$u+w$, the expansions of
$Y(a,u+w)Y(b,w)c$ and $Y(Y(a,u)b,w)c$ agree.

Specify an invertible $\mathsf R[u,u^{-1}]$-linear exchange $\mathcal S(u)$
on $(V\otimes_{\mathsf R}V)[u,u^{-1}]$,
Laurent polynomial on each tensor input. It fixes
both vacuum factors.

Assume both shift equations
\[
 [T\otimes1,\mathcal S(u)]=-\partial_u\mathcal S(u),
 \qquad [1\otimes T,\mathcal S(u)]=\partial_u\mathcal S(u),
\]
the Yang--Baxter identity, and the following field-exchange
identity. For each $a,b$, require an integer $N\geq0$
such that $(z_1-z_2)^N\mathcal S(z_1-z_2)(a\otimes b)$
has no pole and, for every $c$,
\[
 (z_1-z_2)^N Y(b,z_2)Y(a,z_1)c
 =(z_1-z_2)^N Y(z_1)(1\otimes Y(z_2))
       \bigl(\mathcal S(z_1-z_2)(a\otimes b)\otimes c\bigr).
\]
The two sides in this equality use their respective
ordered formal expansions. Also assume the hexagon
\begin{equation}\label{cbc:rc:field-hexagon}
 \mathcal S(z)(Y(w)\otimes1)
  =(Y(w)\otimes1)\mathcal S_{23}(z)\mathcal S_{13}(z+w),
\end{equation}
where the last factor is expanded in nonnegative
powers of $w$. We leave unitarity as an explicit
condition to be tested below.

Use the partition coefficient modules of
\eqref{cbc:ct:partition}, with all normal derivatives
retained, external tensors over $\mathsf R$, and blocks
ordered by their largest labels.
Write $r^\circ_{ij}$ for the principal part of
the field of block $i$ at block $j$, leaving every
spectator state in the output. For ordered distinct
blocks define
\begin{equation}\label{cbc:rc:exchanged-residue}
 X_{ab}=\mathcal S_{ba}(z_b-z_a),\qquad
 r_{ij}=r^\circ_{ij}X_{i,j-1}\cdots X_{i,i+1}
       \quad(i<j).
\end{equation}
The shift equations make $X_{ab}$ horizontal.
The exchange factors are regular in a collision
coordinate whenever their spectator separation
is invertible. Their Taylor expansions therefore
enter the finite principal-part rule
\eqref{cbc:ct:spectator}.

\begin{proposition}\label{cbc:rc:vacuum-exchange}
Put $u=z_1-z_3$ and $\delta_u=[u^{-1}]$ in
the normal principal-part module. On the open
where the second position is distinct from the
first and third,
\begin{equation}\label{cbc:rc:exchange-detector}
 r_{13}(u^{-1};a,b,\mathbf1)
   =\delta_u\,\mathcal S_{21}(z_2-z_3)(a\otimes b).
\end{equation}
The output is written with the merged $13$ state
first and the spectator second. The canonical
flip restores the root order $2|13$.
The full system of residues determines both
$\mathcal S$ and the fields.
\end{proposition}
\begin{proof}
By definition $r_{13}=r^\circ_{13}X_{12}$.
Creation at the last input gives $e^{uT}a$.
The exchange is regular as a series in $u$
near zero on this spectator open. The principal
part of $u^{-1}$ times that regular series is
its constant coefficient times $\delta_u$.
This is~\eqref{cbc:rc:exchange-detector}.

The map placing a tensor in normal degree zero
is injective, so it recovers $\mathcal S$.
For every integer $m$, the coefficient of
$\delta_u$ in $r_{12}(u^{-m-1};a,b)$ is the
coefficient of $u^m$ in $Y(a,u)b$.
Thus the binary principal parts recover the
entire Laurent field, including its regular
coefficients.
\end{proof}

This use of the vacuum is stronger than
remembering an empty-configuration coaugmentation.
It fixes a state and its creation map in a
nonempty coefficient, with the spectator still
present. Applying a scalar functional to that
spectator would change the reconstruction data.

\begin{example}\label{cbc:rc:omitted-exchange}
Let $V=\mathsf R[x_1,x_2,\ldots]$, $J=x_1$,
$Tx_n=nx_{n+1}$, and use the Wick fields
of~\eqref{cbc:ct:kernel} with $\kappa,\lambda\in\mathsf R$ and
$P(u)=\kappa u^{-2}+\lambda u^{-1}$.
The same finite contraction formula defines
these fields even when $P$ is not even.
Its exchange is
$\mathcal S(u)=\exp(C_{21}(-u)-C(u))$.
Each exponential terminates on polynomial
inputs. We compute the defect of omitting
this exchange from the outer residue.

Take states $(J,J,\mathbf1)$ and multiplier
$s^{-1}$, where $s=z_1-z_3$, $t=z_2-z_3$.
Set $u=s-t$. The first consecutive principal
part is that of $P(u)/(u+t)$:
\[
 \kappa u^{-2}t^{-1}+\lambda u^{-1}t^{-1}
                         -\kappa u^{-1}t^{-2}.
\]

The full normal transfer to the common small
diagonal is
\[
 u^{-a}t^{-b}\longmapsto
 \sum_{m=0}^{b-1}\binom{a+m-1}{m}s^{-a-m}t^{-b+m}.
\]
It follows by differentiating the simple
double principal part, since differentiation
in $t$ at fixed $u$ becomes $\partial_s+\partial_t$.

The two consecutive channels and the outer
one without exchange are consequently
\[
 L=-\frac\kappa{st^2}+\frac\lambda{st},
 \qquad R_{\mathrm{ct}}=0,
 \qquad N^\circ=\frac\kappa{st^2}+\frac\lambda{st}.
\]
Each multiplies the vacuum. The second channel
vanishes because creation is regular in $t$.
Thus $L-R_{\mathrm{ct}}+N^\circ=2\lambda\mathbf1/(st)$.

The exchange $X_{12}$ contributes
$2\lambda(s-t)^{-1}\mathbf1\otimes\mathbf1$.
With the multiplier $s^{-1}$, its two principal
parts contribute $-2\lambda\mathbf1/(st)$,
which cancels the defect. The uncorrected defect is nonzero for every $\lambda\ne0$,
also over nonreduced $\mathsf R$. Indeed $2$ is invertible,
and the vacuum is a free direct summand.
\end{example}

\begin{theorem}[Two vacuum partners]
\label{cbc:rc:two-vacua}
For the field and exchange data above, the
commutator of the two crossed residues is
\begin{equation}\label{cbc:rc:crossed-defect}
 r_{24}r_{13}-r_{13}r_{24}
   =r^\circ_{24}r^\circ_{13}X_{23}(X_{21}X_{12}-1).
\end{equation}
Put $p=z_1-z_3$, $q=z_2-z_4$, and $d=z_3-z_4$.
On $d\ne0$, its value on two vacuum partners is
\begin{align}
 &(r_{24}r_{13}-r_{13}r_{24})
       (p^{-1}q^{-1};a,b,\mathbf1,\mathbf1)\notag\\
 &\qquad=[p^{-1}][q^{-1}]
       (\mathcal S_{12}(d)\mathcal S_{21}(-d)-1)(a\otimes b).
 \label{cbc:rc:two-vacuum-defect}
\end{align}
The following conditions are equivalent: unitarity
$\mathcal S_{21}(-z)\mathcal S(z)=1$, commutation of
the crossed residues on the full four-input
coefficient, and vanishing of
\eqref{cbc:rc:two-vacuum-defect} for every $a,b$.
The equivalence holds over nonreduced $\mathsf R$.
\end{theorem}
\begin{proof}
The hexagon~\eqref{cbc:rc:field-hexagon}, followed
by principal part in the collision coordinate,
gives
\[
 X_{2,(13)}r^\circ_{13}
       =r^\circ_{13}X_{23}X_{21}.
\]
The coefficients are expanded before taking
the principal part, as in
\eqref{cbc:ct:spectator}.
Thus $r_{24}r_{13}=r^\circ_{24}r^\circ_{13}
X_{23}X_{21}X_{12}$.
For the opposite order, blocks $1,3$ are
adjacent after merging $2,4$, so
$r_{13}r_{24}=r^\circ_{13}r^\circ_{24}X_{23}$.
The two plain residues act on disjoint
state pairs and commute. Subtraction proves
\eqref{cbc:rc:crossed-defect} without unitarity.

After $X_{21}X_{12}-1$ acts on
$a\otimes b\otimes\mathbf1\otimes\mathbf1$,
the third factor is still the vacuum, so
$X_{23}$ acts identically. The remaining
operator is
\[
 \mathcal S_{12}(d+p-q)\mathcal S_{21}(-d-p+q)-1.
\]
It is regular in $p,q$ near zero when $d$
is invertible. Creation at the two partners
adds $e^{pT}$ and $e^{qT}$. The double
principal part of $p^{-1}q^{-1}$ times
this regular series is its constant
coefficient, giving~\eqref{cbc:rc:two-vacuum-defect}.

Unitarity implies commutation by
\eqref{cbc:rc:crossed-defect}, and commutation
implies the stated vacuum vanishing.
Conversely the normal-degree-zero inclusion
\[
 (V\otimes_{\mathsf R}V)[d,d^{-1}]
 \hookrightarrow
 (V\otimes_{\mathsf R}V)[d,d^{-1}]
        [\partial_p,\partial_q][p^{-1}][q^{-1}]
\]
is injective. Vacuum vanishing therefore
gives $\mathcal S_{12}(d)\mathcal S_{21}(-d)=1$.
Invertibility reverses the factors to give
unitarity. The argument neither specializes
to residue fields nor cancels a field
evaluation map.
\end{proof}

When the exchanged residues satisfy the relations
of Proposition~\ref{cbc:ct:relations}, the finite residue
transfer in Chapter~\ref{cbc:chap:transfer} applies to them.
Its higher products use
the collision of the first and last blocks,
after the exchanges through all intervening
blocks. The vacuum detectors show why those
blocks and the chosen exchange cannot be
discarded during a comparison.

\begin{problem}\label{cbc:rc:full-reconstruction}
Construct the global coefficient category and
comparison functors relating the ordered word
bar--cobar construction to the full collision
atlas and to covariant field systems.
Specify the exact collision functors, the
coefficient dual and its topology, and the
vacuum maps on nonempty configurations.
A field realization must recover all Laurent coefficients through
Proposition~\ref{cbc:rc:vacuum-exchange}. It must also recover the
exchanged outer residues, both vacuum insertions and the compatible
maps for nested collisions.

For the Heisenberg case, compare the native
transverse cochain monoid and its marked
endomorphism system with the completed
coefficient centers of this chapter.
Construct the comparison on full mapping
spaces, specifying the admitted actor
category. The nonrepresentability theorem
excludes all pointed factorizing coalgebras
as a representing domain for $\Phi_C$.
Finality among transverse monoid actions,
with their multiplication coherences and
monoid morphisms, remains to be decided.
\end{problem}
